\documentclass[../main.tex]{subfiles}
\begin{document}
%==========================================TIÊU ĐỀ TẬP========================================
\begin{table}[h]
  \begin{adjustwidth}{-1cm}{}
    \begin{tabular}{>{\centering\arraybackslash}p{6cm}|>{\raggedright\arraybackslash}p{12.5cm}}
      \multirow{5}{6cm}
      % Cột bên trái: ÉPISODE và số tập
      {\\[-40pt]
        \flaregothic\fontsize{20pt}{20pt}\selectfont \centering
        {\contour{errorfive}{\textcolor{black}{ÉPISODE}}}\color{black}\\
        \burbank\fontsize{72pt}{72pt}\selectfont
        \includegraphics[height=3.12359cm]{../../../../Image/Book?/number/num28.png}
        \\[18pt]
        % \flaregothic\fontsize{14pt}{14pt}\selectfont
        % \color{epattr}BIENVENUE, \uline{\color{Banpire} Mel} !
        % \flaregothic\fontsize{18pt}{18pt}\selectfont
        % \color{epattr}ÉPISODE PERDU
        \color{black}
      }
       &
      % Dòng thứ nhất: tên tập tiếng Nhật
      {\fotskip\fontsize{18pt}{18pt}\selectfont \ruby{週}{しゅう}\ruby{末}{まつ}の\ruby{楽}{たの}しみ}
      \\[6pt]
      % Dòng thứ hai: tên tập tiếng Anh
       & {\cafeteria\fontsize{18pt}{18pt}\selectfont Lively Weekend Days}                                                                                                                                                                             \\[4pt]
      % Dòng thứ ba: tên tập tiếng Việt
       & {\vnmsans\fontsize{18pt}{18pt}\selectfont Ngày cuối tuần hoạt náo}                                                                                                                                                                           \\[8pt]
      % Dòng thứ tư: Nhân vật xuất hiện
       & {\caviar{\fontsize{14pt}{14pt}\selectfont Nhân vật: \emp{korone} \emp{okayu} \emp{fubuki} \emp{mio} \emp{subaru} \emp{matsuri} \emp{suisei} \emp{botan} \emp{sora} \emp{polka} \emp{flare} \emp{luna} \emp{kanata} \emp{towa} \emp{roboco}}}
      % \\[6pt]
      % Dòng cuối cùng: Thời gian diễn ra của tập (thường sẽ là ngày phát hành)
      % & {\continuum\fontsize{16pt}{16pt}\selectfont Ngày phát hành: 11/06/2022}\\[8pt]
    \end{tabular}
  \end{adjustwidth}
\end{table}
%===========================================NỘI DUNG==========================================
\color{black}

\pov{Hikawa Kuon}

Tuần đầu tiên ở Aogami đã khép lại. Trong đầu tôi, mọi khoảnh khắc vẫn quay cuồng như đoạn phim quay chậm không dứt.

Aogami, nơi từng là đống đổ nát ngổn ngang từ lời nguyền các vòng lặp, nay đã trỗi dậy như chưa từng chịu đau thương. Người dân hồi sinh. Ký ức được khâu vá. Không gian như bị xóa trắng và viết lại từ trang đầu. Còn tôi, Hikawa Kuon, lẻ loi giữa lớp 1-C, nam sinh duy nhất.

\dots Phải chăng có kẻ nào khác rơi vào cảnh ``độc nam giữa đàn thiên nga'' như tôi? Cả lớp toàn nữ, mỗi gương mặt đều bản sao của idol \hll\ thuộc thế giới cũ. Kể ra ngoài, chắc bị chê hoang tưởng. Nhưng Game-san bảo: đây là vũ trụ \hll\ ERROR, nên việc bạn học giống holomem là chuyện thường ngày.

\dots Ban đầu tôi cũng ráng tự thì thầm vậy.

Những ngày tiếp theo trôi qua y hệt thời cấp ba ở Etsunan: đến trường, học bài, học thêm, về nhà\dots\ đơn điệu như hơi thở.

Ngoại trừ những chuyện không thể coi như chuyện gió bay được.

Tiếng khóc vang vọng cuối hành lang. Bóng nhỏ kimono rách lướt ngoài khung cửa. Lời van xin vọng lên từ dưới đất.

Chúng chưa gây hại\dots\ chỉ là \emph{``chưa''} thôi, chứ không phải là \emph{``không bao giờ''}.

Điều lạ là tôi không phải là người duy nhất nhìn thấy.

Kaigou-chan, Suzuki.

Shino-san, Sakura-san, Yuka-san.

Thậm chí là cả Yuko-sensei, bằng trực giác nhạy như thú săn mồi của cô. Cũng là người mà tôi cũng không muốn bị dây dưa vào nhất.

Sáu người chúng tôi giống như bị đặt vào chung một bàn cờ, buộc phải nhìn thấy những quân cờ vô hình đang di chuyển từng nước đi mà không ai trong chúng tôi nắm bắt được. Yuko-sensei thì đứng ngoài lề, nhưng ánh nhìn cô thì thầm: cô đã thấy quá nhiều.

Nhiệm vụ lúc này: chỉ lặng quan sát.

Không vội vàng.

Không để lộ manh mối.

Không cho vòng lặp một lý do quay trở lại.

Và trong lúc đó, vẫn phải sống như một học sinh bình thường. Một cuộc đời học đường khác, trong một thế giới bị viết lại.

\dots Nhưng nói đi cũng phải nói lại.

Việc trở lại ghế nhà trường cấp ba không tồi như tôi từng mường tượng. Tôi đã quen sống theo nhịp nơi mà những hạn chót cứ treo lơ lửng trên đầu, nhồi nhét cả tuần vào lịch làm việc dày đặc, chẳng để lại một khe hở nghỉ ngơi. Tôi từng tin ``trưởng thành'' đồng nghĩa với cày cuốc suốt đêm, tích góp từng đồng, dọn đường cho tương lai. Các lễ hội, bữa tiệc đón chào, sự kiện kỷ niệm của đất nước đối với tôi chỉ là thứ xa xỉ phẩm.

A-san gợi ý tôi nên dừng lại. Tanigo-san cũng khuyên can. Thậm chí bố mẹ tôi và Ryuuhou cũng bày tỏ quan ngại.  Nhưng tôi vẫn lắc đầu. Tiền - nói thẳng, nói thật - là thứ tôi thèm muốn. Giảm việc đồng nghĩa với giảm thu nhập, tôi không chấp nhận điều đó.

Để rồi giờ đây, khi ngoái đầu ngẫm lại, tôi giật mình thấy mình chỉ còn công việc dài như bất tận và những lần gục xuống bàn vì làm việc quá sức.

Vì vậy, khoảng lặng này có thể là hơi thở hiếm hoi để tôi chậm lại. Vừa tìm lối thoát khỏi Aogami, vừa nếm lại mùi vị học trò từng vô tình bỏ lỡ.

Có thể tôi cũng cần một lần soi gương để tìm lại mình.

Aogami, thị trấn nửa mơ nửa thực.

Nếu đây là kịch bản có sẵn, tôi sẽ bước lên sân khấu.

Nhưng lần này, luật chơi do chúng tôi định đoạt.

\ct

Tôi quay lại với hiện tại khi ánh sáng buổi sáng tràn qua khe rèm cửa sổ.

Thứ bảy, tuần đầu tiên của chúng tôi tại nơi này.

Với học sinh, đó là ngày nghỉ. Với tôi, một đứa thuộc ``câu lạc bộ về nhà'' chính hiệu, đây lại càng là ngày nghỉ đúng nghĩa. Ít nhất là cho tới tuần sau, khi tôi bắt đầu đi làm thêm ngoài giờ học. Còn bây giờ, không bài vở, không lịch trình bắt buộc, không phải nghĩ xem hôm nay có ca nào ở công ty hay không (thực ra là có, nhưng nó có ý nghĩa gì khi tôi bị mắc kẹt ở đây chứ?).

Trong bữa tiệc hôm qua, Yuko-sensei khuyến khích chúng tôi tận dụng cuối tuần đi khám phá thị trấn. Tôi gật đầu qua loa rồi\dots\ thôi. Thú thật, sau mấy ngày tiếp xúc đám đông và chuỗi phiêu lưu thoát chết kéo dài, tôi chỉ muốn một ngày hiếm hoi chẳng làm gì: ở nhà, cắm mặt vào game, ăn gì cũng được - y hệt hồi tôi vẫn còn là một thanh niên lười chảy thây nằm nhà trước khi nhận việc.

Kaigou-chan và Suzuki lại khác. Họ thích giao lưu, cười đùa, kết thêm bạn. Cuối tuần của họ là thế. Còn tôi, chỉ cần góc yên tĩnh là đủ.

Sáng thứ bảy ở nhà Hikawa\dots\ cũng vì thế mà không được yên như mong đợi.

Tôi mở mắt khi đồng hồ điểm bảy giờ, muộn hơn thường lệ hai tiếng.

Vừa lăn người kéo chăn định ``nướng'' thêm thì, *RẦM RẦM* tủ quần áo bị kéo, đóng sầm, móc treo lạch cạch.

``\vie{\dots Ồn quá\dots}''

Tôi lồm cồm ngồi dậy, dụi mắt. Suzuki mải lựa đồ trước tủ, mặt căng như dây đàn, không hay tôi đã thức. Trong lúc đó, cô chị thì dựa tường, cầm điện thoại, thỉnh thoảng liếc xem em chọn xong chưa.

``\vie{Đi đâu thì đi,}'' tôi ngái ngủ càu nhàu, ``{cần gì ồn ào thế\dots}''

Suzuki giật mình quay lại: ``\vie{O-Onii-sama dậy rồi ạ?}''

Cô chị nhếch mép: ``\vie{Tưởng cậu ngủ tới trưa chứ.}''

``\vie{Tôi cũng nghĩ vậy,}'' tôi đáp, ``\vie{nếu hai người không lập cà lập cập như thế.}''

Không ai bận cãi, tôi cũng chẳng buồn hỏi thêm. Nhìn qua biết rất rõ là họ mải chuẩn bị cho buổi đi chơi, chẳng ý định lôi tôi theo. Càng tốt.

Nếu tôi nhớ không nhầm, tôi có nghe thoang thoảng rằng từ bữa tiệc hôm qua, hai chị em nhà nó được các bạn học rủ đi ra ngoài thị trấn. Tôi không chắc liệu họ có mời tôi không, nhưng nếu có thì tôi cũng chẳng bận tâm lắm.

*BÍNH BOONG* *BÍNH BOONG*

Tiếng chuông cửa ngân vang hai lần, ngắt quãng bởi khoảng lặng đủ để tôi nhận ra đây không phải ảo giác buổi sáng.

Từ bên ngoài cửa sổ, tôi nghe thấy những tiếng cười nói ríu rít, như thể có bàn tay vô hình xếp đặt sẵn khung cảnh. Tôi kéo vội chiếc áo khoác mỏng, bước ra sảnh, ngáp dở rồi mở cửa.

``Ồ, các cậu đến tập hợp sớm thế?'' giọng tôi còn đậm mùi buồn ngủ, không giấu nổi sự thờ ơ có chủ đích.

Nanase-san đứng trước, mái tóc xanh khẽ rung theo gió giữa xuân tháng Tư; phía sau cô, Shiina-san khoanh tay nhìn tôi với nét mặt hầm hố như mọi lần, còn Honoka-san nhịp nhịp gót giày thể thao.

Shino-san giương nụ cười mỏng, đưa cao xấp vé màu cam.  ``Bọn mình định ra ngoài thị trấn,'' cô nói, giọng có chút tươi tỉnh hơn mọi ngày (tôi nghĩ bởi đây là lần đầu cô thực sự được bạn rủ đi chơi). ``Vé gần hết rồi, nhóm mình còn dư đúng một chiếc thôi.''

``Không phải tôi ki bo kẹt xỉn gì đâu,'' cô cháu gái của hiệu trưởng trường nói, toát lên vẻ thanh tao của một gia đình quyền quý nhưng gần gũi. ``Thật sự thì rạp chỉ còn có một vé thôi. Rất xin lỗi các cậu vì sự việc này.''

``Vậy à\dots\ Tiếc thật nhỉ,'' tôi buột miệng cười chát, không biết nên buồn hay nên vui. ``Nghe có vẻ như bộ phim đó đắt khách nhỉ.''

Nanase-san nhẹ nhàng nói thêm, giọng vẫn dịu nhưng có chút hồi hộp lạ thường: ``Đúng là như thế. Tôi cũng đã định rủ nhà cậu với Shino-san hồi đầu tuần rồi. Shino-san nghe xong thì gật đầu ngay, còn Shiina-san và Honoka-san thì\dots''

``Tôi chỉ đi theo vì bị Honoka ép,'' Shiina-san cắt ngang, mặt vẫn hầm hố nhưng tai đã đỏ ửng. ``Cô ấy nói nếu không đi thì sẽ post ảnh ngủ dạo lên bảng tin lớp đấy.''

Honoka-san cười toe toét, vỗ vai Shiina-san: ``Cậu mà không đi thì ai còn lắng nghe cậu cằn nhằn về plot hole của phim chứ? Đây là hy sinh vì nghệ thuật đấy!''

``Có mà rảnh háng không có gì để làm thì có,'' cô bạn tóc xám bật lại. ``Lúc nào cũng oang oang cái mồm lên, nhức tai khó chịu.''

Tôi nhìn Nanase-san, chợt nhớ lại: trong một vòng lặp đã bị xóa, cô cũng từng mời tôi và Shino-san như vậy, cũng ngay trong ngày đầu tiên chúng tôi nhập lớp. Giọng điệu lịch sự đến mức lạnh lùng, như thể việc tôi có đi hay không cũng chẳng ảnh hưởng gì đến cô cả.

Còn bây giờ, cô đang ở đây, trước mặt tôi, với một lời tương tự, nhưng có sự đồng thuận của Shino-san, có sự miễn cưỡng của Shiina-san, và có sự nhiệt tình thái quá của Honoka-san. Một lời mời không còn là tình cờ, mà là có chủ đích.

Tôi không biết nên cảm thấy thế nào. Trong vũ trụ mà tôi đang cố gắng sống như một học sinh bình thường, đây là lần đầu tiên tôi cảm thấy mình được nhớ đến từ trước, chứ không phải là một lựa chọn cuối cùng khi mọi người đã đủ vé.

Vừa dứt câu, phía sau họ còn thêm một tốp nữa: Hotaru-san dẫn đầu, theo sau là bộ đôi Kanade-san và Kana-san, cùng với Yuki-san và Mitsuki-san xếp hàng ngay ngắn. Cô nàng tóc hồng vẫy tay, đôi mắt cười híp lại. ``Nhóm mình cũng thiếu đúng một chỗ nè!'' Kana toe toét, giơ ngón tay như thể đang chụp ảnh quảng cáo.

Trong thế giới hoàn hảo ấy, vòng lặp mà Shino-san còn mải mê với những ngày tháng bất tận không đau thương, Hotaru-san từng cất cao giọng rủ rê: ``Shino-chan, cậu về cùng nhóm mình sau giờ học đi!'' Lúc đó cô ấy chỉ cười mờ mịt, gật đầu theo phản xạ, đôi mắt cô ấy như đang nhìn xuyên qua họ, hướng về một kịch bản đã được viết sẵn. Tôi, cùng với Kaigou-chan và Suzuki, những người đứng ngoài lúc đó, cảm thấy mình thừa thãi đến lạnh người. Không ai nhìn tôi. Lời mời gọi đó nghe giả tạo tới mức chẳng khác nào như một câu lệnh được lập trình sẵn và được lặp đi lặp lại.

Nhưng giờ đây, nó lại trở nên chân thực hơn bao giờ hết. Trong lời mời gọi này, tôi cảm nhận được sự chân thành của Hotaru-san, của Yuki-san, và của cặp đôi của câu lạc bộ Phát thanh. Nó làm cho tôi cảm giác họ thật sự sống, chứ không còn như những búp bê vô hồn lặp đi lặp lại.

Mitsuki-san bước lên một bước, giọng không giấu nổi sự nũng nịu như ở thế giới thực: ``Thật ra nhóm tui cũng định mời Shino-san, nhưng nhóm Nanase đã nhanh tay hơn. Tức ghê\~{}'' Cô bĩu môi nhẹ, không giận, chỉ hơi tiếc.

Shino-san gãi má, cười trừ. ``Xin lỗi nhé, lần này mình đã hứa trước rồi. Thôi thì, chúng ta hẹn lần sau nhé.''

Tôi liếc qua liếc lại hai hàng người, rồi nhìn Kaigou-chan và Suzuki đang lùi về phía cầu thang, cả hai đã kịp thay đồ chỉn chu, đồ đạc cũng được lên gọn ghẽ. Mọi thứ quá rõ ràng: họ cần một ``nam chính'' để lấp chỗ trống duy nhất, nhưng không ai dám nói thẳng.

Tôi hít một hơi thật sâu, cảm nhận mùi bánh mì nướng còn vương trong nhà, rồi phẩy tay như xua đi một đám ruồi cam chịu. ``Nếu đã không đủ chỗ thì đành chịu thôi. Tôi ở nhà cũng được. Hai người cứ tận hưởng đi.''

Không cần cân nhắc lâu, cũng chẳng phải giả vờ khiêm tốn. Tôi thực sự muốn một ngày không tiếng ồn, không bài tập, không lịch trình, chỉ có một ngày yên tĩnh, thích làm gì thì làm, hoàn toàn tự do.

Suzuki khựng lại nửa nhịp, đôi mắt hai màu thoáng mất nét hồn nhiên, chỏm tóc ngố trên đỉnh rũ xuống. Nỗi hụt hẫng hiện rồi biến nhanh như tia chớp, kịp để cô khoác lên vẻ hiểu chuyện. ``V-Vậy ạ\dots\ Onii-sama nghỉ ngơi cũng được\dots'' giọng cô nhỏ xíu, như sợ tiếng nói sẽ vỡ nếu to thêm một xíu.

Kaigou-chan thì im lặng lâu hơn. Cô đưa tay vuốt lọn tóc mới làm, ánh mắt xanh nhìn thẳng vào tôi, không trách móc, chỉ là thăm dò. Tôi biết cô đã đoán trước kết cục này: tôi không phải kiểu người thích chen chân giữa đám đông, càng không dễ bị lung lay bởi lời mời nửa vời. Dẫu vậy, trong tầm nhìn ấy vẫn lấp lánh chút nuối tiếc, như thể cô mong tôi đổi ý, dù cả hai đều hiểu điều đó chẳng khác nào bắt mèo leo cây thông.

Cuối cùng, cô ấy gật đầu, khẽ mỉm cười, nụ cười dành cho cậu con trai cứng đầu. ``Ừ. Ở nhà thì nhớ khóa cửa cẩn thận đấy. Quà sáng tôi để trên bàn rồi.''

``Này, chờ chút!'' Yuki-san vỗ tay, thu hút mọi ánh nhìn. ``Giờ cả hai nhóm đều chỉ còn đúng hai chỗ trống. Nhóm Nanase thiếu một, nhóm bọn mình cũng thiếu một. Ai về đội nào đây?''

Suzuki giơ tay trước tiên, mắt sáng lấp lánh. ``Em muốn sang nhóm Mitsuki-san! Mình cũng đang tò mò lắm!'' cô bé quay sang Kaigou, nhỏ giọng, ``Onee-sama, em đi nhé?''

Kaigou-chan cắn nhẹ môi, đưa mắt nhìn Nanase-san, rồi lại nhìn vé cam trên tay Shino-san. ``Thật ra\dots\ chị chưa từng vào rạp bao giờ. Nên\dots\ chị cũng muốn thử xem phim.'' Cô khẽ nắm tay Suzuki, ``Nhưng để em một mình thì liệu có hơi\dots''

Hotaru-san liền ôm vai cô em gái của cô, cười tươi. ``Yên tâm, tôi sẽ trông chừng 'em bé' này cẩn thận, không để lạc mất một sợi tóc nào đâu!''

Suzuki đỏ bừng mặt, giật mình. ``Mình có phải em bé đâu mà\dots'' cô bé lầu bầu, nhưng vẫn dí sát vào cô nàng da sẫm đó, giấu nụ cười khe khẽ.

Nanase-san bất giác lùi nửa bước, ánh mắt thoáng nặng. Cô khẽ lắc đầu, như muốn xua đi hình ảnh mơ hồ: một cái tát giữa hành lang tối, tiếng khóc nghẹn, và cảm giác bị bỏ rơi mà cô không thể nhớ rõ.

Shiina-san đặt tay lên vai cô, thì thầm đủ nghe, ``Có tôi ở đây rồi, cậu lo gì.''

Kaigou-chan cũng gật nhẹ, giọng đều đều, ``Tôi chỉ muốn xem phim thôi, không có ý gì đâu.'' Nhưng đôi mắt xanh của cô vẫn dõi theo bàn tay cô nàng tóc xanh run rẩy, như vừa nhìn thấy ma.

``Đ-Được rồi. Tôi tin cậu,'' cô cháu gái nói.

Hai nhóm lần lượt quay gót. Tiếng bước chân lộp cộp trên bậc tam cấp rồi xa dần, tiếng cười nói nhỏ lại, rồi biến mất sau cánh cửa đóng kín.

Tôi dựa vào khung cửa, nhìn theo bóng Kaigou-chan và Suzuki khuất dần ra xa. Bảy năm trôi qua kể từ ngày bố mẹ họ về nơi chín suối, hai chị em họ sống dựa vào tiền trợ cấp và những buổi làm thêm khuya muộn. Lúc đó rạp chiếu, hay trung tâm thương mại là thứ xa xỉ không cần thiết. Vậy mà hôm nay, Kaigou-chan nắm vé cam như người vừa trúng xổ số, còn Suzuki thì nhảy chân sáo lên xe đạp của Hotaru-san. Có lẽ, với họ, cuộc đi chơi này không chỉ là xem phim, đó là lần đầu tiên được chọn, được nhớ tên, được mời như bao người bình thường khác.

Tôi thì khác. Tôi đã từng đi hết mọi rạp chiếu trong thành phố, đã từng ngồi hàng ghế đầu với đôi tai bị ù nặng vì âm thanh quá lớn, và cũng đã từng đi qua những trung tâm thương mại lớn ở Kawachi cùng với bố mẹ vào những ngày nhà thiếu đồ dùng. Một ngày nghỉ đúng nghĩa với tôi giờ đây chỉ cần bốn bức tường và chiếc TV, không cần phải dấn thân vào những chỗ đông nghẹt người làm gì.

Không khí trở lại mức yên tĩnh nghe được tiếng tủ lạnh gầm gừ. Tôi đóng cửa, thở phào, một cảm giác hỗn độn giữa thoát tục và chút ngậm ngùi khó tả. Có lẽ, ngày mai tôi sẽ hối hận, nhưng hôm nay, tôi chọn làm kẻ ở lại, người giữ lấy nhịp thở chậm giữa thị trấn Aogami vẫn còn ngổn ngang bí ẩn.

\ct

Ban đầu, tôi chỉ định lau dọn nhà cửa sơ sơ để bớt lộn xộn. Chẳng cần dọn dẹp cầu kỳ, cứ phủi bụi nhẹ, gom sách vở tung tóe lại, gấp chiếc chăn còn vắt vẻo trên giường. Những thao tác nhỏ ấy, thực hiện trong đầu trống rỗng, bất ngờ mang lại cảm giác thư thái lạ thường.

Về bữa sáng mà Kaigou-chan đã chuẩn bị sẵn cho tôi. Mở tủ lạnh kiểm tra, ngoài chiếc bánh mì nướng đơn giản ấy, bên trong gần như trống trơn. Chỉ còn vài món ăn vặt và vài lon nước ngọt.

``\vie{\dots Chiều tối đi chợ vậy.}'' Tôi đóng tủ lạnh lại, vừa nói vừa gật gù với chính mình.

Sau khi dọn dẹp căn nhà tạm gọn, tôi ngồi xuống giường, đặt ngón tay lên mặt đồng hồ đeo tay.

``\eng{Game,}'' tôi thì thầm, ``\eng{tình hình ra sao rồi?}''

Mất vài nhịp tim, giọng nói Pháp quen thuộc mới vang lên, trầm ấm và nằm sát ranh giới tiềm thức.

``\eng{Tôi vẫn ở đây.}''

Tôi không cần hỏi vì sao ông im hơi lặng tiếng mấy ngày qua, câu trả lời đã rõ.

``\eng{Cậu hiểu lý do mà,}'' Game tiếp, như đọc thấu suy nghĩ. ``\eng{Tôi chỉ cất tiếng khi không có người ngoài, hoặc khi những kẻ đã `thấy' tôi xuất hiện.}''

Tôi khẽ hừ một tiếng. ``\eng{Vậy ông cũng để ý lũ nhóc đó.}''

``\eng{Thấy,}'' ông ta đáp, ``\eng{và nghe.}''

Tôi chỉnh thẳng lưng. ``\eng{Ông đánh giá sao?}''

``\eng{Chúng không có ác ý,}'' ông nói chậm. ``\eng{Những linh hồn ấy\dots\ chưa trọn vẹn. Chúng hiện hình nhiều hơn khi có âm thanh, sinh hoạt đông người, hay thứ gì gợi `niềm vui'. Bữa tiệc hôm qua là minh chứng.}''

Tôi nhớ lại ánh đèn lập loè, những đốm sáng lơ lửng Suzuki kể. ``\eng{Thế\dots\ ý ông là chúng bị mắc kẹt?}''

``\eng{Gần đúng,}'' Game đáp. ``\eng{Chúng tìm cách muốn bắt chuyện. Một thứ mong muốn gì đó được để ý. Và để\dots\ được giải thoát.}''

Tôi lặng đi mấy giây. ``\eng{Còn lối ra?}''

``\eng{\dots Chưa,}'' ông ta thẳng thừng. ``\eng{Chưa có manh mối rõ ràng. Nhưng tôi sẽ tiếp tục.}''

``\eng{Ừ,}'' tôi nói. ``\eng{Cứ từ từ.}''

Cuộc đối thoại kết thúc. Tôi đứng dậy, duỗi cánh tay, rồi trở lại với kế hoạch ban đầu.

Sau khi dọn dẹp xong, tôi lần mò vào bếp, moi ra gói snack còn sót lại trên kệ cao nhất. Tôi bật TV, không thèm chọn kênh, để mặc màn hình nhảy qua những chương trình quảng cáo rực rỡ rồi thả mình xuống giường.

Tôi thở ra một hơi dài, cảm giác như vừa tháo được cục tạ trên vai.

``\vie{Có lẽ\dots\ cứ như thế này cũng đủ rồi. Chắc hôm nay mình sẽ thư giãn một ngày, không làm một cái gì cả.}''

Chưa kịp bóc gói, tiếng chuông cửa cắt ngang khoảnh khắc yên bình.

*BÍNH BOONG*

Tôi nheo mắt, giữ nguyên tư thế. Tôi sẽ không làm gì cả.

*BÍNH BOONG* *BÍNH BOONG* *CỘC CỘC*

Tiếng chuông trở nên sốt ruột, thêm mấy cú gõ cửa dồn dập, không dịu dàng, không do dự, kiểu người sẽ đứng đó đến trưa nếu cần.

Tôi thở dài, ném gói đồ ăn vặt lên bàn, lẩm bẩm trong miệng: ``\vie{Chắc lại bọn trẻ trâu rảnh rang không có gì làm đây mà\dots}''

Tôi lê bước ra hành lang, tay vừa chạm nắm cửa đã thấy mình như diễn viên bước vào cảnh quay không thoát được. Mở ra, và khựng lại.

Trước mặt tôi là cả một đội hình nữ sinh đúng nghĩa ``hội đồng cố vấn'', không ai khác ngoài những bạn học cùng lớp.

``\dots Cái lề gì thốn.''

Akane-san đứng nghiêng, nụ cười toe toét như đã đoán trước mọi phản kháng của tôi. Yuka-san nheo mắt, đôi tai mèo phụ kiện khẽ rung như thể ngửi được mùi lười biếng trong nhà. Phía sau, Inari-san bưng túi giấy to oạch, Hanabi-san vẫy vẫy tay như đi đại hội. Kaoru-san đang kéo lê Miku-san tới, một tay nắm cổ tay cô, tay kia vẫy vẫy. Cô nàng bờm sói mếu máo, bước đi lê bước, mắt nhìn xuống đất như kẻ bị dẫn độ, nhìn vừa thương vừa buồn cười.

``Yo,'' Kaoru-san lên tiếng trước, giơ tay chào như đang đứng sân bóng. ``Tới chơi đây.''

Yuka-san nghiêng đầu, mái tóc tím khẽ rủ xuống trán. ``Kuon-kun ở nhà thật kìa\~{}''

Hanabi-san liếc nhanh vào trong, khóe môi nhếch. ``Nghe nói hôm nay cậu định ở ẩn à?''

Tôi không đáp. Thay vào đó, tôi từ từ\dots\ đóng cửa lại, như kéo rèm kết thúc vở kịch sớm.

Chưa kịp khép hẳn, bỗng một bàn tay chặn ngang, lực đẩy nhẹ nhưng đủ khiến cánh cửa giật lại.

Tôi nhìn thẳng vào người vừa chặn cửa, thở dài chán chường: ``\dots Các cậu có hẹn trước không?''

``Không,'' Hanabi-san đáp tỉnh queo. ``Tự dưng thấy thích thì đến thôi. Dù gì cũng là hàng xóm mà.''

``Tôi thì không có gì để làm, lại còn được Yuka và Akane mời nên tới thôi,'' Kaoru-san tiếp lời.

Tôi thở dài, tay áp lên trán, khẽ lắc đầu như không tin nổi với những lý do đầy ngẫu hứng đó. ``\dots Thế thì xin mời đi về,'' tôi nói, giọng đều đều như con robot hết pin, ``tôi không có kế hoạch tiếp khách.''

``Ê!'' Akane-san vờ nhăn mặt, khẽ phồng má lên. ``Bất lịch sự thế\~{} Khách tới chơi mà lại đuổi đi thế hả!''

Yuka-san cười khẽ, giọng mềm như bông gòn. ``Chỉ vào ngồi một lát thôi mà.''

Hanabi-san gật phụ họa, nhấc túi giấy lên quá đầu. ``Bọn mình mang đồ ăn theo đó. Có bánh dorayaki nóng, trà hoa cúc tự nấu, cả bịch khoai tây chiết tay mà Wata-san làm nữa này\~{}''

``Mình biết là cậu không phải là người thích giao tiếp,'' Inari-san tiếp lời, hai tay đan lại với nhau, ``nhưng\dots\ ít nhất mình muốn chúng ta gắn kết với nhau hơn.''

Tôi liếc vào túi giấy, rồi nhìn lại căn nhà trống trải phía sau. Phía sau lưng, bàn ăn trơ trụi, sofa không vỏ gối, ti-vi đang chiếu quảng cáo nước rửa bát. Mùi snack phô mai thoang thoảng như lời cám dỗ cuối cùng.

``\dots Không được sao?'' nàng bờm chó ngước lên tôi, đôi mắt long lanh như chó con bị bỏ rơi xin được nhận nuôi. Trời ạ.

Sau ba giây cân não (thực ra là tôi thấy mình đã thua trước khi bắt đầu), tôi buông tay khỏi cửa, xoay người bước vào.

``\dots Vào đi,'' tôi thở dài, như tháo van. ``\vie{Mịa, coi như đi tong luôn ngày yên tĩnh.}''

Kaoru-san cười tươi rói, đập vai Miku-san. ``Thấy chưa, biết ngay là Kuon-kun sẽ đầu hàng mà.''

Cô nàng bờm Sói đen hé miệng, giọng nhỏ như muỗi kêu: ``\dots Tôi bảo là để tôi ở nhà mà\dots'' Cô cúi đầu, má đỏ ửng, như kẻ vừa bị bắt cóc tới đây. ``Mà rõ ràng trông cậu ta còn như thể không muốn tiếp kia kìa.''

``Rồi mọi thứ sẽ quen thôi\~{}'' Hanabi-san nói, kéo mọi người vào trong.

Tôi để mặc cho cả bọn tràn vào phòng khách, mỗi người một câu, một giọng, khiến căn nhà vốn yên ắng bỗng chốc náo nhiệt hẳn lên.

Ít ra\dots\ tôi tự an ủi mình: ``\textit{Còn hơn phải vác xác ra chỗ đông người ngoài kia.}'' Thôi thì, đằng nào họ cũng đã cất công tới, đuổi đi cũng không nỡ, ở nhà một mình cũng chẳng có gì làm ngoài ăn với ngủ cả. Thêm vài bạn học trong phòng cũng không mất gì, đúng chứ?

Kaoru-san là người chiếm ghế trước, gác chân rất tự nhiên như thể đây là nhà cô ấy. Hanabi-san thì tò mò nhìn quanh mấy kệ sách, còn ``bộ tứ GAMERS'' bày đồ ăn vặt ra bàn. Tôi đứng dựa vào khung cửa một lúc, nhìn cảnh tượng ấy mà có cảm giác lạ lẫm.

Một nhóm học sinh tới chơi nhà vào ngày cuối tuần. Cảnh này, với tôi, tưởng chừng đã thuộc về một cuộc đời khác. Hồi tôi sống ở thế giới gốc, điều này không bao giờ xảy ra. Chỉ đơn giản là vì tôi không muốn mời họ thôi, ai cũng có công việc riêng của mình.

Nhưng không. Điều tưởng chừng bất khả ấy lại đang diễn ra trước mắt tôi, ngay tại lúc này.

Tôi quay mặt về phía cửa sổ.

Bên ngoài, con phố Aogami yên ả đến mức gần như không thật. Ánh nắng trưa rơi đều trên hàng cây, không có tiếng còi xe dồn dập, không có lịch làm việc dày đặc, không có chuỗi tin nhắn thúc giục. Chỉ là một ngày nghỉ bình thường.

``\textit{Trông chẳng khác nào những ngày chủ nhật đối phó với mấy bài vở của các gái ở văn phòng ghê\dots}'' tôi nghĩ thầm, khóe miệng nhếch lên rất nhẹ.

``Kuon-kun?'' Yuka-san quay sang. ``Cậu không ngồi xuống à?''

``\dots À.'' Tôi lững thững đi tới, ngồi xuống đầu bàn.

Và thế là, một ngày nghỉ của tôi cùng với các cô gái lớp 1-C cũng bắt đầu từ đây.

\ct

``Căn nhà này rộng hơn tôi tưởng nhỉ,'' cô nàng bờm Chó thốt lên, đôi mắt không giấu nổi sự ngạc nhiên khi ngó nghiêng khắp gian phòng khách rộng thênh thang. ``Không giống kiểu nhà thuê của học sinh chút nào.''

``Eh, bình thường thôi,'' tôi đáp gọn lỏn, không muốn dây dưa vào chi tiết. Chỉ có điều, đến căn nhà của tôi sống ở Kawachi cũng chỉ rộng đâu đó bằng một góc của căn phòng này, nên việc được sở hữu căn phòng rộng này đã là một điều khiến tôi ngạc nhiên và có gì đó lạ lẫm lúc đầu. Nhưng rồi mọi thứ cũng quen dần, đâu lại vào đấy.

Hanabi-san nghiêng đầu, mái tóc đuôi ngựa khẽ rung theo động tác. ``Bọn tôi tự nhiên ngồi nhờ thế này, có khiến cậu khó chịu không?'' Nụ cười cô vẫn rạng rỡ như mọi khi, nhưng phía sau đó là chút e dè khó tả.

``Khách đã đặt chân tới cửa thì phiền cũng hóa không,'' tôi nói, giọng khô như ngói nung, nhưng giờ không hề có ý đuổi. ``Các cậu dự định cắm trại ở đây tới bao giờ?''

``Tới khi cậu phải dùng chổi đuổi bọn tôi ra ngoài ấy,'' Kaoru-san cười ha hả, vỗ vỗ bụng như thể đang kể chuyện tiếu lâm.

``Và tôi vừa làm thế thật,'' tôi chống hông, chưng hửng. ``Nhưng nhìn cách các cậu thê này thì có dọa ma cũng chẳng thèm về.''

Inari đỏ bừng hai gò má, lí nhí: ``C-Chỉ là qua chơi thôi\dots\ không có ý quấy rầy thật mà.''

Tôi lần lượt đảo mắt qua từng người: Akane-san với vẻ tinh nghịch vừa chuẩn bị trêu chọc, Yuka-san với ánh nhìn xảo quyệt như mèo vừa trộm được cá, Inari-san có phần rút rè như thỏ non, Hanabi-san tràn trề năng lượng như pin sạc đầy, Miku-san lặng lẽ đứng cuối hàng, còn Kaoru-san thì\dots\ đúng chất Suba-chan ở thế giới gốc, sẵn sàng xông pha bất cứ đâu.

Không khí họ mang tới nhẹ tênh, không chút mùi tang tóc nặng trĩu. Không vòng lặp nào đang quay, không kịch bản nào được kích hoạt. Đơn thuần chỉ là lũ bạn cùng lớp rủ nhau qua nhà chơi cuối tuần.

``\dots Tùy mọi người,'' tôi kết luận sau vài giây cân não. ``Chỉ cần đừng đập vỡ cái gì là được. Nhà vừa mới dọn đấy.''

Cả bọn đồng loạt cười khúc khích. Âm thanh trong trẻo ấy lan tỏa, khiến căn phòng khách vốn lạnh lẽo bỗng ấm lên vài độ, như có luồng gió xuân lướt qua khe cửa sổ.

Ngồi đã ấm chỗ, cô nàng tóc ngắn lên tiếng: ``Ở nhà cậu có trò gì chơi không? Chứ ngồi nói chuyện không hoài cũng chán.''

Tôi đáp ngay: ``Nhà có vài bộ boardgame. Để đâu rồi nhỉ\dots'' Nói đoạn, tôi đứng dậy, lục trong tủ dưới kệ vô tuyến, kéo ra mấy hộp cũ. Kaoru-san và Hanabi-san lập tức sáng mắt, còn Inari-san thì cẩn thận đỡ giúp tôi bày lên bàn.

Tôi thực ra cũng không nhớ mình từng mua những hộp này. Hôm qua, khi dọn dẹp lại căn nhà sau khi đã dỡ đồ, tôi tình cờ thấy chúng nằm gọn gàng trong ngăn kéo dưới gầm tủ, như thể ai đó đã cất cẩn thận để dành cho một dịp như hôm nay. Có lẽ là đồ của người thuê trước, hoặc của chủ nhà để lại, nhưng dù sao thì chúng cũng đang ở đây, và có vẻ còn nguyên vẹn.

``Cái này chơi thế nào?'' Miku-san hỏi nhỏ, chỉ vào chiếc hộp bàn cờ.

``Luật đơn giản thôi,'' tôi nói, vừa giải thích vừa sắp xếp quân cờ. ``Ai về đích trước thì thắng.''

Ván đầu tiên trôi qua trong tiếng cười, đặc biệt là khi Akane-san liên tục đi sai luật còn Kaoru-san thì cố tình ``chơi xấu'' để trêu. Yuka-san cười khẽ mỗi lần tôi nhắc lại luật, còn Hanabi-san thì hào hứng tới mức suýt hất đổ cả bàn.

Giữa lúc cả bọn đang tranh luận ầm ĩ về một nước đi, tôi vô tình đá phải thứ gì đó cứng cứng dưới chân.

*CẠCH*

Tôi cúi xuống.

Một chiếc máy chơi game thế hệ cũ nằm lẫn trong góc kệ, phủ một lớp bụi mỏng. Trông giống như của nhà Nin, cùng thời với trò mà nhân vật chính là một thợ sửa ống nước.

``\dots Cái này là gì vậy?'' tôi khẽ nhíu mày.

Kaoru-san tò mò thò đầu xuống. ``Máy game à? Của cậu hả?''

``\dots Không nhớ là có,'' tôi đáp thành thật.

Yuka-san khoanh tay, nghiêng đầu. ``Có khi là đồ cũ trong nhà từ trước?''

Tôi cầm nó lên, phủi bụi. Dây nối còn nguyên, tay cầm vẫn ở đó. Một cảm giác kỳ lạ thoáng qua, như thể có thứ gì đó trong căn nhà này đã ở đây từ rất lâu, nhưng tôi chưa từng để ý tới.

Hanabi-san reo lên: ``Cắm thử đi! Biết đâu còn dùng được!''

Tôi do dự một nhịp, rồi gật đầu. ``Được rồi\dots''

Khi cắm dây vào TV, màn hình lóe lên trong chốc lát trước khi hiện giao diện đơn giản, cổ điển đến mức hoài niệm.

Akane-san há hốc miệng. ``Ôi, đúng kiểu ngày xưa luôn!''

Miku khẽ mỉm cười. ``Trông\dots\ thân thuộc thật.''

Tôi nhìn màn hình, rồi nhìn lại căn phòng đang dần đầy tiếng nói cười.

Một ngôi nhà bình thường. Một buổi cuối tuần bình thường. Bạn bè ngồi quanh bàn, tranh cãi vì trò chơi.

Không vòng lặp. Không cái cảm giác bị theo dõi.

Chỉ là một cuộc sống bình dị. Một ngày cuối tuần êm đềm mà tôi đã mong chờ từ rất lâu.

Và trong khoảnh khắc đó, tôi chợt nhận ra: Có lẽ, nơi này đang thực sự trở thành một ``không gian an toàn'' cho lớp 1-C\dots\ và cho cả chính tôi.

\begin{center}
  \pov{Hikawa Suzuki}
  \uline{Trung tâm thành phố.}
\end{center}

Ba mươi phút kể từ khi nhóm năm người chúng tôi rời nhà, và cũng là lần hiếm hoi tôi tách khỏi Onee-sama, nhưng với một tâm thế khác.

Khi bị kéo vào The Void, tôi như kẻ lạc đường giữa một thế giới đen kịt, mỗi bước chân đều vang lên tiếng vỡ của nỗi sợ. Tách biệt khỏi Onee-sama khi ấy là điềm báo của cô đơn vô tận, là cảm giác bị bỏ rơi giữa không gian trống rỗng, nơi tiếng gọi tên tôi chỉ trả lại tiếng vọng của chính mình. Nhưng bây giờ, trên con đường rợp bóng cây ra ga Aogami, tôi vẫn cảm nhận được hơi ấm của chị như luồng sóng vô hình bám theo từng vòng quay bánh xe đạp.

Cái tách biệt này không gợn chút đen tối; nó như khoảnh trời mở ra để tôi thấy mình là ai khi không cần núp dưới bóng của chị, một cô gái có thể tự cân bằng trên đôi bánh xe, tự quyết định nên dừng lại ngắm hoa anh đào hay đạp thật nhanh qua góc phố quen. Tôi nghe thấy tiếng Hotaru-san cười khan, tiếng Kana-san huýt sáo giai điệu bài hát mới, tiếng chuông xe đạp của Kanade-san réo lên như bản nhạc đồng ca đón tôi vào cuộc sống nơi ánh sáng. Tôi không còn là mảnh vỡ lạc lõng, mà đã là một mắt xích được mời gọi, được cần đến, và khoảng trống Onee-sama để lại không phải hố đen, mà là khoảng trời để tôi tự bay, biết chắc rằng khi tôi mỏi cánh, bàn tay chị vẫn ở đó, sẵn sàng đón lấy.

Chúng tôi cùng nhau đến sảnh trung tâm thương mại gần ga Aogami, cả năm người bước ra khỏi bãi đỗ xe đạp chung một lúc. Tôi vừa gạt chống, vừa nghe Kana-san nói vang: ``Hôm nay phải chơi cho đã!'' còn Kanade-san gật đầu, mắt híp cười; Yuki-san ngân nga giai điệu quen thuộc, trong khi Mitsuki-san đã hướng thẳng về phía dãy cửa hàng, ánh mắt sáng rỡ.

Hotaru-san đi bên tôi, vỗ vai rồi chỉ vào bản đồ treo tường. ``Tụi mình bắt đầu từ tầng hai nhé!'' Cô ấy nói lớn, như thể cả nhóm đã sẵn sàng cùng lúc, không ai đến trước hay đến sau.

``Vậy, kế hoạch thế nào?'' cô nàng tóc tím hỏi, tay chống hông như một MC bắt đầu chương trình.

Hotaru-san lập tức đáp, giọng chắc nịch: ``Ăn trưa trước, rồi đi dạo mua sắm. Nếu còn thời gian thì ta có thể hát karaoke.''

``Có karaoke nữa ư!?'' Yuki-san quay phắt lại, hai mắt long lanh như thể tìm được cơ hội.

Tôi hơi chậm nhịp, tay vô thức siết nhẹ quai túi xách, rồi khẽ hỏi: ``Ờm\dots\ về tiền nong thì sao?''

Cả bốn người nhìn tôi trong một giây, rồi Kana-san cười xòa: ``Bao cậu.''

``Đúng đó na,'' Mitsuki-san gật đầu, giọng dịu dàng. ``Hôm nay bọn mình rủ cậu đi chơi mà.''

``Nhưng mà\dots'' tôi lúng túng, cảm thấy mình đang nhận quá nhiều. ``Mình không muốn\dots\ ăn chùa.'' Dẫu biết rằng xuất thân của hai chị em tôi là một gia đình nghèo, tiết kiệm từng đồng một, thắt chặt chi phí tới mức không dám tiêu gì cho bản thân, nhưng tôi không muốn bị coi là kẻ mang ơn.

Hotaru-san khoanh tay, vẻ nghiêm túc. ``Không phải ăn chùa. Là bọn này muốn vậy đấy, Suzuki.''

Tôi cắn môi. Cảm giác được mọi người đối xử tốt khiến tôi ấm áp, nhưng cũng có chút áy náy. ``Vậy\dots\ cho mình đóng góp một phần nhé. Dù ít thôi cũng được.''

Kana-san nhìn sang Kanade-san, rồi lại nhìn tôi. ``Ừm\dots\ được. Nhưng cậu đừng có ngại ngùng quá nha.''

Kanad-sane khẽ nói, giọng nhỏ nhưng rõ: ``Bọn mình đi cùng nhau là vui rồi. Không cần phải khệ nệ quá đâu.''

Tôi thở nhẹ, mỉm cười. Ít nhất, tôi cũng có thể góp một phần vào ngày hôm nay.

Mitsuki đã là người dẫn đầu. ``Đi nào, khu mua sắm phía trước vui lắm na\~{}''

Cô ấy bước rất tự nhiên giữa dòng người, vẻ ngoài dịu dàng nhưng cách đi thì dứt khoát. Tôi theo sau, mắt đảo quanh để chắc chắn không ai trong nhóm bị lạc, cạnh bên là Hotaru-san đang chỉ trỏ từng cửa hàng dễ thương, thỉnh thoảng dừng lại lâu hơn mức cần thiết trước mấy con thú nhồi bông.

``C-Cái này\dots\ dễ thương thật,'' cô nàng lẩm bẩm, rồi đỏ mặt khi Kana-san quay sang trêu: ``Ủa, ai mới nói không thích mấy thứ này nè?''

``K-Không có!'' cô bạn phản bác ngay, nhưng `ztai đã đỏ bừng.

Cùng lúc đó, Mitsuki-san kéo tôi ghé vào một tiệm phụ kiện. ``Cậu xem thử cái này đi. Hợp với cậu đó\~{}''

Tôi cầm chiếc kẹp tóc trên tay, những ngón tay chạm vào bề mặt mát lạnh của kim loại. Một cảm giác chân thực đến lạ lùng. ``Đ-Đẹp thật\dots'' tôi bối rối đáp.

Ở phía sau, bộ đôi ``Quỷ - Thiên thần'' bắt đầu tranh luận về việc mua quà lưu niệm cho lớp.

``Cái này làm giải thưởng trò chơi được nè,'' Kana nói.

Kanade-san nghiêng đầu: ``Nhưng mà liệu có hơi\dots\ kỳ không? Tôi thấy bà toàn chọn quà kỳ quái không à.''

``Kỳ quặc cái gì, là sáng tạo!''

Hai người cứ thế mà tranh luận qua lại, khiến chúng tôi bật cười. Hai người bọn họ lúc nào cũng vậy, từ bên trong lớp đến ra ngoài thành phố, cứ lúc nào ở bên là lại như thế.

Không khí nhẹ nhõm, ồn ào theo cách dễ chịu. Tôi tự hỏi, liệu thế giới bên ngoài kia có bao giờ rực rỡ và sống động đến thế này không? Hay chỉ khi ở bên họ, mọi thứ mới trở nên sắc nét như vậy?

\ct

Chúng tôi dừng lại ở một gian đồ lưu niệm. Đủ thứ: móc khóa, mũ, khăn, những món nhỏ xinh.

``Ôi, cái mũ này dễ thương nè!'' Kana-san chộp lấy một chiếc mũ hình tai mèo, đội lên đầu rồi quay vòng. ``Thấy sao, Suzuki-san?''

``Dễ thương thật!'' tôi mỉm cười, mắt sáng lên. ``Rất hợp với cậu nữa.''

Cô ấy nghiêng đầu, nheo mắt như mèo vừa nghe được lời ru. ``Thật chứ? Hay chỉ là lịch sự thôi?''

``M-Mình khen thật mà!'' tôi vội vã đáp, hai má ửng hồng.

Kana-san cười khúc khích, tháo mũ ra rồi quay sang Kanade-san đang thử bờm sừng. ``Này, Kanade, cậu cũng thử đi!''

Kanade-san giật mình, lùi lại nửa bước vì vô tình va vào kệ trưng bày. Mấy món đồ rung nhẹ, chực chờ rơi xuống.

``Ơ-Ơ\dots\ xin lỗi\dots!'' cô nàng luống cuống, tay vội đỡ nhưng lại càng làm nó rung thêm.

Tôi phản xạ bước nhanh lên trước, một tay giữ lấy chiếc kệ, tay kia nhanh chóng đỡ lấy mấy hộp quà sắp tuột khỏi mép. ``K-Không sao đâu. Cẩn thận đi chứ.''

Mitsuki-san che miệng cười khẽ. Hotaru-san nhìn tôi, ánh mắt như vừa phát hiện ra điều gì đó: ``Suzuki-san, cậu trông giống ra dáng người chị ghê.''

Tôi hơi sững lại, chỉnh lại mấy món đồ cho ngay ngắn rồi cười nhẹ: ``Mình\dots\ Mình không dám nhận\dots''

Kana quay lại, tháo mũ ra. ``Thấy chưa, có Suzuki-san đứng đây là yên tâm hẳn.''

Tôi cảm thấy má hơi nóng. Giữa gian hàng nhỏ ấy, trong tiếng nói cười của các bạn, tôi chợt nhận ra mình không còn là người đứng ngoài nữa. Không phải là một dữ liệu, không phải là một hình chiếu. Tôi đang ở đây, hít thở bầu không khí này, chạm vào những món đồ này, cùng họ, trong một buổi cuối tuần rất\dots\ bình thường.

Và kỳ lạ thay, cái thực tại này khiến lòng tôi nhẹ đi.

\begin{center}
  \pov{Hikawa Kaigou}
  \uline{Rạp chiếu phim.}
\end{center}

Nhóm chúng tôi đi cùng nhau tới khu rạp chiếu phim ở tầng trên. Không ai bị bỏ lại phía sau, đúng nghĩa là tập kết. Bây giờ đã là chín giờ ba mươi sáng.

Lần đầu tiên kể từ ngày bị kéo vào The Void từ thế giới cũ, tôi rời khỏi Suzu mà không còn cảm giác như bị ai đó rút cạn không khí trong lồng ngực. Tôi vẫn còn nhớ cảm giác bất lực, gào tên em đến khản cổ giữa bóng tối đặc quánh, mỗi tiếng vang chỉ trả lại tiếng vọng của chính mình, một âm thanh trống rỗng đến thảm hại. Giờ đây, khoảng cách giữa hai chị em không còn là vực thẳm nuốt chửng, mà chỉ là vài bước chân đi trước hay đi sau trên hành lang trung tâm thương mại. Tôi vẫn cảm thấy hơi ấm của em như luồng sóng vô hình bám theo từng nhịp tim, nhưng không còn là xiềng xích nữa; đó là bảo bối mềm mại tôi có thể cất vào túi áo, tin chắc rằng khi cần, chỉ cần mở ra là sẽ lại thấy Suzu cười tươi.

Hôm nay, tôi gửi lại em cho nhóm của Kana-san, những người bạn học đang bước cạnh mình mà có thể tin tưởng. Trong lúc đó, tôi có thể tận hưởng cuộc sống của một cô gái bình thường, ở trong một thế giới bình thường, dù chỉ là trong chốc lát.

Tôi đi ở phía sau một chút, vừa đủ để nhìn thấy bốn người kia bước cạnh nhau: Nanase-san đi thẳng lưng, ánh mắt luôn để ý xung quanh; Honoka-san thì ríu rít nói không ngừng; Shiina-san tay đút túi áo, vẻ mặt thản nhiên nhưng rõ ràng là đang quan sát mọi thứ như đang dè chừng một thứ gì đó.

Shino-san thì\dots\ trông khác hẳn những gì tôi từng biết về cô ấy trong những vòng lặp trước đó. Cô ấy cười. Một nụ cười rất tự nhiên của một cô nữ sinh. Giống như mọi ngày.

Cô cháu gái nhà hiệu trưởng khẽ vẫy năm tấm vé trong tay, giọng điềm đạm như thể chỉ hỏi thời tiết. ``Suất mười giờ, phim học đường. Tôi đã giữ chỗ từ hôm qua.''

Honoka-san sáng mắt, reo lên như nhóc được kẹo. ``Kìa, phim đang gây bão trên mạng ấy hả? Tớ còn chưa kịp đọc review của mọi người nữa!''

Cô bạn tóc nâu khẽ giật mình, mắt dính vào tấm vé, giọng nhỏ như sợ gió thổi bay. ``À\dots\ mình có lướt qua trailer rồi. Nghe đồn bộ phim này nói về những điều ước mà học sinh nào cũng từng thầm giữ.''

Shiina-san liếc tôi, khóe môi nhếch đầy ẩn ý. ``Nghe có vẻ sến sẩm thế nào ấy.''

Cô bạn bờm Cáo lập tức bênh vực, hai tay chống hông. ``Sến một chút mới thấm chứ! Shiina-chan đúng là khô cứng ghê! Chẳng có tí thanh xuân tuổi trẻ gì cả!''

``Kệ tôi.''

Tôi gật nhẹ, điệu bộ cảnh giác. ``Được, miễn nam chính đừng có màn tỏ tình sến súa kiểu nhảy cẫng lên bàn ăn là được. Trông nó khó chịu lắm''

Cô bạn bờm Sư tử cười khẽ, đập vai tôi. ``Có thì tôi đạp cho một phát, cậu yên tâm.''

\ct

Cả năm đứa kéo nhau lê bước tới quầy bắp nước như đoàn khách bất đắc dĩ vừa mới thống nhất xong rằng ``xem phim thì phải có đồ ăn thức uống''. Shiina-san như thường lệ, là người đặt chân lên bậc đá trước tiên; cô nàng nghiêng người về phía thủy tinh đựng bắp rang, mắt sáng lên, giọng vẫn thản nhiên đến mức tưởng như đang hỏi thời tiết:

``Mua năm phần có được giảm giá không chị?''

Một cô gái trạc tuổi chúng tôi - có vẻ là nhân viên ở đây - cười trừ, hai tay vẫn cầm cái xẻng bắp run run, như thể đang bị hỏi đòi nợ. Shino-san và Nanase-san, như hai bảo mẫu quen việc, đồng thanh kéo cô bạn lùi lại.

``Shiina-san, thôi mà\dots''

``Đừng gây khó dễ cho người ta.''

Honoka-san đứng kế bên, mắt tròn xoe, khẽ huýt sáo một giai điệu vô thưởng vô phạt như thể đang thưởng thức một vở kịch nhỏ giữa giờ ra chơi.

Tôi lùi thêm nửa bước, khoanh tay dựa vào cột, nhìn cảnh tượng mà không nhịn được cười khẽ: ``Quả nhiên chị đại xứ Aogami vẫn giữ phong độ nhỉ.''

Shiina-san quay nửa thân, liếc xéo: ``Cậu nói gì đó?''

``Không có gì,'' tôi nhún vai, ``chỉ thấy quen quá.''

Cuối cùng, sau ba phút mặc cả vừa đủ để nhân viên không phải gọi bảo vệ, cả bọn cũng ôm về hai túi bắp và bốn cốc nước. Tôi lẳng lặng không đặt món gì, chỉ đứng nép vào cạnh quầy như chiếc bóng dưới ánh đèn neon chập choạng ban sáng.

Shino-san thấy vậy, khẽ nghiêng đầu: ``Sao cậu không mua gì đi?''

Tôi cười gượng, đá chân vào viên gạch: ``Tiếc tiền. Mà, mình ăn chung được mà.''

Nanase-san nghe vậy, không nói gì, chỉ gật đầu với nhân viên rồi lấy thêm một hộp bắp rỗng và một chiếc cốc trống. Cô mở túi bắp to ra, rót một nửa sang hộp mới rồi đẩy ngay phần bắp về phía tôi. Shino-san cũng bắt chước, lấy thêm cốc rỗng, lẩy bẩy chia nước ra làm ba, miệng lẩm nhẩm: ``Một phần của cậu, một phần của tớ, một phần\dots\ cũng của tớ luôn, coi như tớ mời.''

Tôi ôm túi bắp, mùi bơ nóng bốc lên làm cay mắt. Cảm giác ấy\dots\ rất bình thường - thứ ``bình thường'' mà tôi chưa từng được trải nghiệm trong đời: được chia, được nhận, không cần cất tiếng xin.

Tôi chợt nghĩ tới Kuon-kun và Suzu -  không rõ giờ này hai người đang làm gì: có lẽ vẫn ngồi chung chiếc sofa cũ, hoặc đang cãi nhau về một bộ boardgame nào đó, linh cảm của tôi mách bảo như vậy.

Ý nghĩ ấy lướt qua nhanh như sợi chỉ khói, nhẹ tênh, nối những mảnh rời rạc của buổi sáng hôm nay lại thành một bức tranh chỉ mình tôi thấy được.

\begin{center}
  \pov{Hikawa Kuon}
  \uline{Nhà Hikawa.}
\end{center}

``\dots Đã giờ này rồi à.'' Đồng hồ trên nhà đã chỉ về lúc mười một giờ trưa.

Tôi vừa để nhóm con gái ngồi ổn định trong phòng khách thì ánh mắt đã tự động lia sang túi đồ Inari đặt trên bàn. Chiếc túi ni-lông từ cửa hàng tiện lợi nhăn nheo, miệng túi hé mở, lộ ra đủ loại snack đủ màu: bánh quy gói đỏ, khoai tây chiên xanh lá, vài lon nước ngọt lấp lánh dưới ánh đèn phòng khách.

``Chừng đó liệu đủ ăn không vậy?'' tôi hỏi, nghiêng đầu nhìn sơ qua. ``Tận bảy người chúng ta đấy chứ chả ít.''

Inari gật gù, hai tay vẫn ôm túi như bảo vật. ``Ăn lót dạ thì đủ rồi đó, Kuon-kun.''

Miku chống tay lên hông, mắt liếc quanh bàn trống trơn: ``Nhưng mà chẳng phải như thế là hơi--''

``Món gì cũng có hạn sử dụng cả,'' Kaoru cắt ngang, vừa nói vừa xoay lon nước ngọt trong tay như kiểm tra nhãn: ``Ăn nhanh kẻo lỡ hạn.''

Akane nhún vai, nói tỉnh bơ như đang thông báo thời tiết: ``Dù sao cũng đừng hi vọng bọn này xuống bếp, nha.''

Tôi khựng lại nửa nhịp, mắt vẫn dán vào mớ đồ ăn vặt. ``\textit{\dots Vậy là xác định không nấu cơm à.}''

Giọng tôi trầm xuống, như kết luận cho một cuộc đàm phán đã rã đám đầy chóng vánh. Đúng theo dự đoán của tôi. Dù sao tôi cũng chẳng định vào bếp từ đầu.

Tôi thở dài nhẹ, vuốt vội mái tóc rối rồi quay xuống sofa. ``Thế thì tôi ra chợ gom đồ cho bữa chiều luôn. Ai cần gì thì gọi món luôn đi.''

Chưa ai kịp mở miệng, tôi đã lần lượt chỉ từng cô nàng.

``Akane-san chắc chắn sẽ gọi bánh mềm hoặc đồ ngọt.''

Cô nàng bờm chó giật bắn. ``Sao cậu biết!?''

Tôi không đáp, quay sang người bên cạnh. ``Hanabi-san thì snack mặn hoặc gì đó cay xé lưỡi.''

Cô nàng trợn mắt. ``Cậu đọc được suy nghĩ người khác à!?''

Tiếp theo là Inari-san. ``Còn cậu thì là trà sữa hoặc mấy món linh tinh ngoài quầy nước đúng không?''

Cô nàng bờm cáo trắng cười khổ, hai má đỏ ửng. ``Mình vừa định nói\dots''

Ba cô gái nhìn tôi như thấy một con quái vật ngoài hành tinh vậy. Tôi đáp lại họ bằng một cái nhún vai: ``Dễ như bỡn.''

Chưa dừng lại, tôi nói tiếp: ``Yuka-san là cơm nắm.''

Cô bạn đứng hình một giây. ``\dots Hả?''

``Cậu thích món đơn giản, dễ ăn, lại no lâu. Nhìn là biết ngay.'' Mặc dù đó hoàn toàn là phỏng đoán, dựa theo món khoái khẩu của Okayu-chan ở thế giới gốc.

Cô bạn bờm mèo khẽ đỏ má. ``Sao Kuon-kun biết rõ thế\dots''

Tôi quay sang Miku-san. ``Miku-san thì là sô-cô-la đúng không nhỉ?'' Cái này tôi biết chắc chắn. Trong một vòng lặp ở toa tàu, tôi thấy cô ấy cứ khư khư cầm thanh sô-cô-la, nên tôi tin như vậy.

``Ể-- đúng rồi!'' cô nàng bờm sói mắt sáng lên. ``Tôi định mua thật đấy!''

``Với lý do `làm đẹp da' chứ gì.''

``\dots Cậu đáng sợ ghê.''

Cuối cùng, ``Còn Kaoru-san thì hẳn nhiên là hộp sữa.'' Cái này cũng không hề bàn cãi, ở vòng lặp trước cũng vậy.

Cô nàng tóc ngắn tròn xoe mắt. ``S-sao cậu biết luôn!?''

``Cậu uống sữa mỗi khi thấy có cơ hội, lại còn thỉnh thoảng lẩm bẩm muốn cao thêm vài phân nữa rồi còn gì.''

Kaoru-san đỏ bừng tai. ``Đừng nói to thế chứ!!''

Căn phòng lặng đi đúng hai nhịp tim.

Và rồi\dots

``Kuon-kun đọc vị người khác giỏi quá\dots'' Inari-san thì thào.

``Cảm giác như bị soi thấu vậy\dots'' Hanabi-san rùng mình.

Akane-san khoanh tay trước ngực. ``Từ nay chắc không giấu nổi gì trước cậu rồi ha\~{}''

Tôi cười khẽ. ``Linh tính mách bảo thôi. Chẳng có gì đặc biệt đâu.''

Tôi xỏ dép, với lấy chiếc ví mỏng trên bàn. ``Vậy tôi đi đây. Ai muốn đi cùng không?''

Dứt lời, cả nhóm con gái đều nhìn nhau.

Inari-san nghiêng đầu. ``Mình ở lại dọn snack ra cho tiện.''

Hanabi-san lắc mạnh. ``Tôi cũng thấy ngại đi nữa.''

Ba người còn lại thì cùng đồng thanh: ``Bọn tôi cũng ở lại!''

``Vậy còn Yuka-san và Akane-san thì sao?'' tôi nhìn về phía cặp bạn thân.

Yuka-san bước lên trước, mắt sáng rõ lên: ``Tôi đi!''

Akane-san lập tức nhảy dựng khỏi sofa: ``Cho tôi theo với\~{}''

Cô nàng tai mèo nhíu mày, khẽ kéo tay cô bạn: ``Bà ở nhà cũng được mà. Đi với tôi làm gì cho mệt người ra.''

``Không!'' cô bạn tai chó nằng nặc, ôm lấy cánh tay Yuka. ``Tôi nhất quyết phải đi với bà! Tôi không muốn bị bỏ lại đâu!''

Tôi nhìn hai người, khẽ thở dài rồi mỉm cười. Cảnh tượng này thật quen thuộc, hình như tôi đã thấy điều đó ở đâu rồi. ``Được rồi, nếu Akane-san muốn vậy thì cứ để cho cậu ấy đi.''

``Yatta!'' cô nàng tóc nâu reo lên hớn hở, còn các cô gái còn lại thấy vậy cũng chỉ cười mỉm với phản ứng này.

Trước khi cả nhóm còn đang cười đùa với nhau, tôi khựng lại một khắc.

``\textit{Không biết Kaigou-chan và Suzuki giờ này thế nào nhỉ\dots}''

Chắc vẫn đang dạo phố, cười nói rộn ràng. Một ngày nghỉ hiếm hoi, đành để họ tận hưởng.

Tôi mỉm cười nhạt, khép cánh cửa nhẹ nhàng sau lưng.

\begin{center}
  \pov{Hikawa Suzuki}
  \uline{Nhà ăn - Trung tâm thương mại.}
\end{center}

% Ăn trưa

Khu ăn uống ở tầng năm đông đúc hơn tôi tưởng tượng rất nhiều. Tiếng nói cười, tiếng dép lê lê lết trên sàn gạch hòa vào mùi dầu mỡ bốc lên từ các quầy, xộc thẳng vào mũi khiến bụng réo lên thình thình, như thể nó đang phản chủ vì đã bỏ đói nó buổi sáng.

``Được rồi, chia đội đi mua nhanh nào!'' Kana-san vỗ tay một cái chát, giọng vang như MC đang dẫn chương trình. ``Hotaru với Suzuki bên này, Mitsuki với Kanade bên kia, còn Yuki--'

``Tôi sẽ chọn quán đông nhất!'' Yuki-san giơ tay cao, mắt long lanh như vừa phát hiện kẹo mới.

``\textbf{Yuki!! Quay lại đây!}'' cặp bạn thân ấy đồng thanh kêu lên, nhưng cô bạn tóc đen đã chạy biến.

Tôi bật cười khẽ, lách qua dòng người theo Hotaru-san. Cô ấy bước nhanh, mái tóc vàng bay nhẹ, để lại hương dầu gội hoa anh đào thoang thoảng. Đứng trước bảng menu đèn neon, cô nghiêng đầu, ngón tay gõ nhịp lên cằm.

``Suzuki, cậu ăn cay tốt chứ?'' cô hỏi, mắt vẫn dán vào tấm ảnh mì kim chi bốc khói.

``Mình\dots\ tạm ổn\dots'' tôi gãi má, nhớ lại lần ăn thử wasabi để rồi đến chảy nước mắt khi lần đầu ăn sushi hồi ở thế giới trước.

``Vậy cậu lấy phần của cậu đi, tôi gọi cay.'' Cô bạn da ngăm quay sang quầy gọi món bằng giọng trầm ấm, dứt khoát. Nghe cô nói ``Một phần mì kim chi đặc biệt, thêm trứng luộc'', tôi chợt thấy ấm áp như đang đi chợ cùng chị gái, thứ cảm giác xa xỉ mà tôi hiếm khi nào trải nghiệm được.

\ct

Mười phút sau khi chúng tôi gọi món, cả năm chúng tôi tụ lại quanh bàn dài gỗ sáng bóng. Mitsuki-san đẩy khay của cô ấy ra giữa, miệng cười tủm tỉm: ``Mọi người cứ gắp chung cho vui na\~{} Còn nhiều đồ lắm á\~{}'' Trên khay là những viên takoyaki còn bốc khói nghi ngút, bên dưới lớp bột vàng ươm là miếng bạch tuộc giòn sật.

Kanade-san ngồi cạnh cô bạn thân, cẩn thận xếp đũa song song, như sợ chúng cọ xát nhau sẽ trầy. Kana-san thì vừa xì xụp mì vừa bình luận: ``Góc trái kia bán gà rán ngon ghê, nhưng hơi đắt nha--''

``Kana, ăn đi rồi nói tiếp,'' cô bạn tóc trắng khẽ kéo tay cô bạn, mặt đỏ ửng vì bị ``bóc phốt''.

Trong khi đó, Yuki-san ngồi đối diện, miệng ngậm ống hút, thân hình lắc lư theo nhịp nhạc chỉ cô mới nghe thấy. Mỗi lần cô lắc mạnh, lọn tóc lại quật vào vai cô nàng tóc hồng, khiến cả bàn cười rộ.

Tôi nhìn quanh, mắt dừng lại trên làn khói bốc lên từ tô mì, trên khuôn mặt rạng rỡ của Hotaru-san đang thổi nóng, trên bàn tay ai đó vô tình chạm vào khi gắp miếng thịt. Không khí ồn ào, đời thường, ấm áp đến lạ. Tôi chợt nhận ra: đây là lần đầu tiên tôi được ngồi giữa một nhóm bạn, ăn một bữa không cần lo lặp lại, không cần đếm thời gian, cũng không cần giả vờ. Chỉ đơn giản là một thế giới bình thường, sống như những con người bình thường, và làm những điều mà người bình thường làm.

Ở thế giới cũ, bàn ăn của tôi (thực chất là một cái mâm) chỉ có một chỗ đối diện Onee-sama; tiếng đũa va vào nhau nghe thật to trong căn bếp vắng, và mỗi câu nói của tôi đều phải cân nhắc kẻo làm chị lo. Khi ấy, tôi từng nghĩ rằng ``bạn bè'' là khái niệm dành cho người khác, những kẻ không bị đánh dấu bởi dòng chữ ``gây nguy hại cho cô em gái bé bỏng'' - là tôi - trên trán. Giờ đây, giữa tiếng cười cợt, tiếng chê mì cay của Kana-san và tiếng ``cẩn thận bỏng'' của Kanade-san, tôi mới hiểu thế nào là được đón nhận. Không còn phải nơm nớp lo sợ việc bị Onee-sama gây khó dễ hay bị bắt nạt nữa.

Tôi bất giác mỉm cười, cảm thấy hơi nóng từ tô mì lan tỏa lên đôi má, lan vào trái tim đang đập chậm rãi, đều đặn. Một buổi trưa rất bình thường, nhưng với tôi, đó đã là cả một kỳ tích: từng bước chân thoát khỏi bóng người chị, từng tiếng nói cười thay cho tiếng vọng cô đơn, từng hơi ấm này đang viết lại định nghĩa ``nhà'' trong tôi.

\ct

Khi khay trước mặt tôi chỉ còn vài giọt nước tương lấp lánh và đồ ăn cũng đã cạn vơi đi gần hết, Yuki-san đập bàn đứng phắt dậy, mái tóc đen đỏ loang như tia lửa. ``Được rồi! Giờ tới phần chính của ngày hôm nay!''

Tiếng ghế kéo rít lên, Kana-san nhíu mày. ``Còn trò gì nữa?''

Tôi còn chưa kịp nuốt miếng mì cuối, cô ấy đã chỉ thẳng qua ô cửa kính, mắt sáng như đèn pha. ``Karaoke! Quán kia mới khai trương, mình đi hiến kỹ năng thanh nhạc cho đời đi!''

Kanade-san giật mình, đũa rơi khỏi tay. ``Sớm vậy!? Giờ nghỉ trưa mà--''

``Càng phải đi chứ!'' cô bạn xoay người, váy tung bay như vũ công rối. ``Ca sĩ nào cũng phải rèn giọng sau bữa trưa, đó là thời điểm thanh quản nở nhất!''

Tôi bật cười khẽ. Cô ấy lại mang cả lý thuyết âm nhạc vào buổi đi chơi. Nhớ hôm qua, Yuki-san cũng hát song ca với Azuki-san tại tiệc chào mừng, giọng ca treo cao như chim sơn ca nhưng lại nhảy loạn xạ như chim bị điện giật.

Kana-san thở dài, lắc đầu. ``Không ai đăng ký khóa học của cậu hết đó\dots\ Cứ xấn xa xấn xổ kiểu thế chẳng ai theo kịp đâu--''

Cô bạn tóc đen không đợi ai phản đối thêm, nắm cổ tay Mitsuki-san kéo đi. ``Đi mà đi mà!''

Mitsuki-san bị giật tung, chỉ kịp cười trừ với chúng tôi. ``Ừm\dots\ thôi thì, đi một chút cũng được. Dẫu sao ta cũng chẳng còn gì khác để làm na.''

Hotaru-san nhún vai, mắt lim dim như mèo vừa ngủ dậy. ``Coi như tiêu bớt năng lượng sau khi ăn vậy.''

Cô bạn tóc tím thở dài, nhìn cô ấy với sự ngán ngẩm như đã đoán định từ trước. ``Thấy chưa, bị kéo đi luôn rồi. Tôi còn lạ gì với cái tính nhanh nhảu của cậu ta nữa.''

Cô bạn thiên sứ gãi má, đỏ ửng. ``Ừ\dots\ nhưng chắc cũng vui.''

Tôi siết nhẹ vạt áo. Karaoke - một từ ngữ xa lạ trong thế giới cũ, nơi tôi chỉ biết ngân nga một mình trong căn phòng trống, giọng hòa vào tiếng vọng của chính mình. Giờ đây, nó lại đến như lời mời gọi từ thế giới ngoài kia.

Hotaru-san quay lại, nắm vai tôi. ``Đi chứ, Suzuki? Hôm nay là ngày nghỉ mà, cứ chơi tới bến đi!''

Tôi nắm lấy bàn tay cô ấy - ấm và dịu dàng đến lạ, rồi gật đầu nhẹ.

Cả nhóm rủ nhau đứng dậy, tiếng cười nói vang lên như hợp xướng rộn ràng. Tôi bước theo họ, tim đập chầm chậm, nhưng không còn là nhịp cô đơn. Một buổi trưa bình thường, đang dần khắc thành ký ức có tên: ttôi, người được mời hát.

\begin{center}
  \pov{Hikawa Kaigou}
  \uline{Rạp chiếu phim.}
\end{center}

Sau khi mỗi người đã có cho mình một phần đồ ăn kèm, cả nhóm bắt đầu lũ lượt bước vào phòng chiếu. Tiếng giày dép lách cách trên sàn gạch sáng bóng, hòa vào tiếng nhạc nền nhẹ nhàng phát ra từ loa treo tường. Những giai điệu piano du dương như đang ru ta vào một thế giới khác.

``Nhanh lên nào, phim sắp bắt đầu chiếu rồi đấy,'' Nanase-san khẽ khàng nói, đi lên đầu. Cô ấy bước chậm rãi nhưng đầy quyết đoán, ánh mắt lướt qua từng hàng ghế như thể đang kiểm tra một bản đồ chiến thuật.  Bốn chúng tôi khẩn trương theo sau.

Cửa phòng chiếu khép lại sau lưng bọn tôi bằng một tiếng *cạch* trầm đục; một thứ âm thanh khô khan, dứt khoát, không khác gì lời tuyên bố: ``Từ giờ, mọi thứ đã bắt đầu đổi khác. Không thể chùn bước thêm nữa.''

Bên trong tối hơn hẳn hành lang bên ngoài, chỉ còn ánh đèn dẫn lối màu xanh nhạt chạy dọc sàn, ánh đèn nê-ôn sáng tựa như những viên đá phát sáng nhỏ xíu, lấp lánh như sao rơi.

Mùi bắp rang bơ lẫn với mùa ghế nệm cũ đặc trưng của rạp phim, một thứ mùi ẩm ướt, pha lẫn hơi ngai ngái của vải bạc và bụi thời gian, khiến tôi có cảm giác rất đời thường.

Một buổi chiều thứ bảy, trong một rạp chiếu phim bình thường, cùng những người bạn đang bước đi bên tôi.

Cô nàng tóc xanh dừng lại ở hàng ghế thứ năm từ dưới lên.  ``Ở chỗ này là đẹp nhất.''   Giọng cô ấy nhỏ nhưng đầy thuyết phục, như thể đây không chỉ là chọn ghế, mà là chọn một vị trí trong cuộc đời: ở giữa, một khoảng cách vừa đủ, không ở phía trước để bị giật mình, nhưng cũng không quá lùi phía sau để bị bỏ lỡ.

Honoka-san theo sau, reo lên khe khẽ:  ``Háo hức quá đi!'' Cô ấy nhún nhảy từng bước, hai tay ôm bắp rang như ôm một con thú bông, mắt long lanh như đèn LED.

Shiina-san ngồi bên cạnh cô bạn tai cáo kia, khoanh tay, liếc nhìn màn hình chưa sáng hẳn. ``Hy vọng không sến sẩm quá.''  Môi cô ấy mím chặt, một thói quen khi đang chờ đợi được trải nghiệm một cái gì đó mới mẻẻ.

Shino-san vào sau cùng, im lặng hơn mọi khi, nhưng đôi mắt cô ấy ánh lên vẻ háo hức rất rõ, như hồ nước phẳng lặng bỗng có tia nắng lọt xuống, rọi thẳng vào đáy.

Tôi ở ngoài cùng, sát lối đi, từ thói quen cũ  hồi còn hay đề phòng đủ thứ. Vẫn còn cảm giác phải sẵn sàng chạy, phải để lối thoát sau lưng. Nhưng hôm nay, tôi không còn muốn chạy nữa.

Tôi chỉ muốn tận hưởng một cuộc sống bình thường của một nữ sinh bình thường, một điều rất xa xỉ mà thôi.

``Thưa quý khách, xin vui lòng tắt nguồn điện thoại hoặc chuyển sang chế độ im lặng. Xin cảm ơn quý khách đã đến với rạp chiếu phim Aogami Cinema. Bộ phim sẽ bắt đầu sau ít phút.''

Giọng nữ ngọt ngào vang lên từ loa trần, như một lời nhắc nhở nhẹ nhàng: ``Hãy tạm biệt thế giới bên ngoài. Bây giờ, chỉ còn chúng ta và bóng tối.''

Ánh đèn dần tắt. Từng dải ánh sáng xanh nhạt lụi tắt như những vì sao bị gió thổi bay. Không gian chìm trong bóng tối đặc quánh, một thứ bóng tối không đáng sợ, mà ấm áp tới kỳ lạ. Như một cái ô khổng lồ, ôm trọn lấy chúng tôi.

Màn hình sáng lên. Ánh sáng trắng xanh loang ra, chiếu rọi lên từng khuôn mặt.

Tôi liếc sang Shino-san. Cô ấy ngồi thẳng đăng, hai tay đặt lên đùi, nhưng đôi mắt mở to đến nỗi có thể đựng cả bầu trời trong đó.

Bộ phim mở đầu bằng khung cảnh sân trường mùa xuân. Những cánh hoa anh đào rơi, học sinh cười nói, một nam sinh vụng về làm rơi cặp trước mặt một cô gái. Khung cảnh quen thuộc đến nỗi như đang nhìn vào một ký ức không thuộc về mình, nhưng sao lại quen thuộc tới lạ.

Honoka-san che miệng cười khúc khích: ``Trời ơi, mở đầu dễ thương ghê!'' Cô ấy nghiêng đầu, tựa vào vai Nanase-san, một cử chỉ tự nhiên đến nỗi khiến tôi chợt thấy hơi ghen tị. Ghen vì sự dễ dàng ấy, cái cách mà cô ấy có thể chạm vào người khác mà không sợ bị đẩy ra.

Cô nàng tóc xanh nghiêng đầu chăm chú theo dõi, ánh mắt nghiêm túc như đang xem tài liệu nghiên cứu.
Cô ấy không cười, không nhúc nhích, mà chỉ quan sát, như thể đang học cách làm người từ những thước phim. Quả đúng là con gái của một gia đình danh giá có khác.

Shiina-san thì hừ nhẹ, có gì đó thô ráp hơn: ``Thằng này đúng kiểu ngốc nghếch.'' Nhưng khóe môi cô ấy hơi nhếch một nụ cười không muốn thừa nhận.

Tôi bật cười khẽ. Chúng tôi đang ngồi đây, cùng nhau, trong một rạp chiếu phim, cười vì một câu chuyện không liên quan đến mình.

Nhưng câu chuyện trên màn ảnh bỗng chốc lắng sâu, lặng lẽ tuôn ra những gam màu trầm. Hoa anh đào không còn rơi, tiếng cười ngây ngô cũng tắt lịm, nhường chỗ cho những buổi học trò ở lại trường muộn: ánh đèn vàng vọt hắt bóng người trên tường trắng, tiếng gõ bàn đều đặn của người trực nhật, và những khoảnh khắc va chạm lỡ làng, tin nhắn gửi nhầm, bài tập đánh rơi, ánh mắt vừa chạm đã vội quay đi như kẻ phạm tội. Cảm xúc vụng về, muốn nói mà không dám, muốn giữ mà chẳng biết cách, len lỏi trong từng khung hình.

Shino-san ngồi bên tôi, lưng thẳng, hai tay đặt lên đùi. Thỉnh thoảng vai cô khẽ run, hòa nhịp với thở của phim—nhanh khi cảnh giật mình, chậm khi màn hình chìm vào im lặng. Tôi cảm giác cô không chỉ đang xem, mà đang sống lại từng mảnh ký ức vỡ vụng của chính mình.

Rồi nữ chính thì thầm, không khóc, không lớn tiếng, như tự nhủ với bản thân năm mười lăm tuổi:

\begin{center}
  ``Nếu được sống lại từ đầu\dots\ tớ muốn làm mọi thứ khác đi. Muốn cười nhiều hơn. Muốn dũng cảm hơn.''
\end{center}

Tôi thấy rõ bàn tay Shino-san siết chặt tay áo, đốt ngón trắng bệch vì dùng lực, như níu kéo điều gì sắp tuột khỏi tầm với. Đôi mắt xanh thẫm long lanh; không phải nước mắt, chỉ là ánh màn hình phản chiếu. Nhưng đủ để tôi hiểu: cô đang khóc bên trong, tiếng khóc không thành lời, chỉ là hơi thở run nhẹ, như chiếc lá rơi xuống mặt hồ, không gợn sóng, nhưng đã từng rung động.

Tim tôi thắt lại. Tôi cũng từng nghe câu ấy, trong đầu, trong mơ, trong vòng lặp bất tận, trong ác mộng lặp lại. Giờ đây, tôi ngồi đây, bên người cũng đang giữ lấy điều chưa dám làm thuở xưa: khát khao tuổi học trò bị ngoại cảnh bóp nghẹt, nay đã gần hơn bao giờ hết.

Tôi không nhìn cô. Chỉ khẽ dịch tay, để lại khoảng trống vừa đủ cho một bàn tay nếu cô muốn chạm.

Trong bóng tối ấm và ánh sáng mờ ảo từ màn ảnh, tôi chợt tự hỏi: ``Mong ước của Shino-san được tiếp tục sống những giây phút đẹp nhất đã thỏa chưa? Hay cô còn điều gì chưa thể thốt thành lời?''

Dẫu vậy, có một điều tôi dám khẳng định, ít ra lúc này: chúng tôi đang thực sự sống, làm những gì mình muốn, và được chính mình kiểm soát.

Nhưng rồi, sự yên bình ấy chẳng kéo dài được bao lâu — nó chỉ như một lớp sương mỏng phủ trên mặt hồ, chờ một cơn gió vô tình để vỡ tan thành gợn sóng cong cong.

Giữa lúc bộ phim chuyển sang một cảnh tối hơn — hai nhân vật chính ngồi dưới mái hiên tránh mưa — ánh sáng trong rạp dịu xuống, và theo phản xạ quen thuộc, tôi liếc nhanh ra phía sau.

Hàng ghế cuối.

Khi nãy vào, tôi chắc mẩm: ``\textit{Chắc nó cũng trơ trụi như những hàng ghế khác thôi.}''

Không một ai ngồi đó cả.

Ban đầu tôi nghĩ đó chỉ là ảo giác. Vậy mà bây giờ\dots\ ở đó có \emph{một cái bóng}.

Không phải hình người rõ ràng, mà giống như một mảng tối đặc hơn bóng tối xung quanh, hơi nhòe viền, như thể không thuộc về thế giới này. Phần trên phình ra như cái đầu cúi xuống, phần dưới co lại — nhỏ hơn nhiều so với người lớn.

Một dáng ngồi. Gối ôm sát ngực. Tựa như một đứa trẻ đang thu mình lại trong sợ hãi.

``\dots!''

Tim tôi khựng một nhịp.

Cái bóng ấy\dots\ đang nhìn về phía trước.

Chính xác hơn, nhìn về phía Shino-san.

Đúng lúc đó, màn hình phim sáng lên. Cảnh trường học ngập nắng, nhạc vui vẻ vang lên, chói hơn bầu không khí vừa rồi.

Trong khoảnh khắc ánh sáng phản chiếu khắp rạp, tôi nhìn rõ hơn.

Không còn bóng mờ nữa.

Ở ghế sau cùng\dots\ một đứa trẻ gầy gò, tóc rối, kimono cũ rách viền, vải sẫm màu như bùn đất thấm nhiều lần.

Nó co gối, hai tay siết chân. Đầu hơi nghiêng. Ánh mắt vô hình, nhưng tôi cảm được: dán vào cô nàng tóc nâu đó.

Honoka-san thì thào: ``Ơ\dots\ phim có đoạn này à?''

Nanase-san cau mày. ``Không ăn nhập gì.''

Shiina liếc sau rồi nhíu sâu. ``Vừa rồi\dots\ có ai ngồi đó phải không?''

Không khí lạnh ngắt. Có vẻ như các bạn của tôi cũng đã nhận được điều gì đó mờ ám. Tuy rằng không có cảm xúc quá mãnh liệt.

Nhưng Shino-san\dots\ Cô rùng mình. Nhẹ, như gió lạnh lướt sống lưng.

Nhưng cô ấy không hét. Không hoảng loạn. Ngược lại, mắt cô mở to, chứa cảm xúc khó tên, mặc dù cô còn không nhìn ra đằng sau, như thể đã đoán định được thứ gì đó còn đeo bám ở nơi này, một ký ức từng chôn sâu chỉ chờ được khai quật lên lần nữa.

Tôi ngoái lần nữa.

Hàng ghế cuối\dots\ bất chợt trống trơn.

Không có bóng hình. Không có đứa trẻ nào. Cũng không còn dấu vết.

Chỉ ánh đèn mờ và phim vẫn chạy, như chưa hề xảy ra chuyện gì.

Tim tôi đập thình thịch. Tôi nắm chặt tay.

Tôi nhìn Shino-san thêm lần nữa. Cô vẫn dán mắt vào màn, nhưng hơi thở chậm hơn, ánh mắt chìm xuống.

Và trong đầu tôi vang rõ: ``\textit{\dots Linh hồn không chỉ lẩn quẩn thị trấn. Chúng\dots\ bám người từng mắc kẹt.}''

Bám theo ký ức. Bám theo những nnỗi đau cũ. Bám theo những vòng lặp đã từng nuốt chửng họ.

\ct

Bộ phim kết thúc trong tiếng nhạc nhẹ và tiếng vỗ tay lác đác của khán giả. Đèn rạp bật sáng, ánh sáng trắng lạnh lẽo lập tức xua tan cái bóng tối ấm áp vừa ôm trọn, đưa chúng tôi về thực tại.

Honoka-san vươn vai, giọng reo vang như chuông gió: ``Phim hay ghê! Tớ suýt khóc đoạn cuối!''

Nanase-san gật đầu hài lòng, tay vuốt phẳng nếp váy: ``Cốt truyện ổn, dù cú twist giữa chừng hơi gượng.''

Shiina-san hừ khẽ, mắt vẫn dán vào màn hình đã tắt: ``Cũng được, không tồi\dots\ nhưng cái đoạn bé gái kimono ngồi sau lưng tụi mình\dots\ tôi cảm thấy kỳ kỳ thế nào\dots''

Cô bạn tai cáo chớp mắt: ``Hả? Cậu nói ai cơ?''

Cô bạn tóc xanh cau mày, đưa tay đỡ trán. ``Cậu nói lảm nhảm gì thế? Bộ mắt cậu có vấn đề gì sao?''

Tôi khẽ lên tiếng, giọng thấp hơn mọi khi: ``Có lẽ là cậu bị quáng gà thật rồi.''

``\dots Cũng có thể.'' Shiina-san nhún vai, tặc lưỡi cho qua.

Shino-san đi bên cạnh mỉm cười, nhưng tôi thấy rõ trong đáy mắt cô ấy vẫn còn vương lại điều gì đó rất sâu, như hồ nước tối có bóng rọi xuống chưa từng được gợn sóng.

Như Kuon-kun đã nói từ trước đó, tôi không muốn càng có nhiều người bị kéo vào chuyện này. Rõ ràng Shino-san cũng biết điều đó, vì thế cô ấy không lên tiếng.

Chúng tôi bước ra hành lang sáng choang, đồng hồ treo tường chỉ 12:34. Ánh đèn neon phản vào gạch bóng loáng, lạnh tanh.

``Mitsuki-san từng nói trung tâm thương mại này có cả khu trò chơi và quán kem nổi tiếng,'' tôi vô thức lên tiếng, đánh lái sang chủ đề khác.

Honoka-san nghe vậy, liền hào hứng vỗ tay: ``Hay đi ăn kem đi! Xem phim xong là phải ăn kem mới đúng bài!''

Nanase-san gật ngay. ``Tôi sẽ gọi vị trà xanh matcha trân châu đen.''

Shiina-san chống hông, vẫn còn nghiêng nửa người về phía rạp như thể muốn quay lại kiểm chứng. ``Miễn đừng kéo dài quá\dots\ Sau đó tôi muốn tỉ thí với mấy cái máy trong khu trò chơi kia.''

Shino khẽ gật đầu, mắt lia tới tôi. Một cái nhìn hỏi ý, nhưng cũng như cầu cứu.

Tôi bật cười nhẹ, cố làm nhẹ bầu không khí: ``Vậy thì đi thôi. Dù sao còn nguyên buổi chiều cho chúng ta nữa.''

Cả nhóm rẽ sang khu ăn vặt. Honoka, Nanase và Shiina dẫn đầu, tiếng cười nói lại vang lên, như thể khoảnh khắc trong rạp phim chưa từng tồn tại.

Tôi cố tình bước chậm. Shino-san hiểu ý, khựng lại bên cửa kính trưng bày mascot bằng nhựa cao su. Dòng người qua lại che khuất hai chúng tôi khỏi tầm nhìn của ba người kia.

Cô thì thầm, giọng mỏng như làn gió: ``\hl{Kaigou-san\dots{} em bé đó\dots{} dường như cũng từng là nạn nhân của nghi thức hiến tế ấy.}''

Tôi quay sang, mắt cô ướt, nhưng không cho phép giọt nước rơi. ``Mình không dám quay lại\dots\ vì sợ em ấy vẫn còn đó\dots\ và sẽ mãi đợi mình nói lời xin lỗi.''

Tôi khẽ đặt tay mình lên tay cô một cách nhẹ nhàng, trấn an: ``Nếu em ấy tìm tôi, tớ sẽ nói hộ cậu\dots\ rằng cậu đã đủ dũng cảm để bước ra khỏi bóng tối. Cậu không cô độc đâu.''

Shino-san siết nhẹ ngón tay, rồi buông, mỉm cười gượng. ``Làm phiền các cậu rồi.''

Cô bước đi trước, tóc nâu khẽ bay. Tôi theo sau, lòng nặng trĩu.

Dù bề ngoài mọi thứ vẫn bình thường\dots\ tôi biết rất rõ: những linh hồn ấy đã bắt đầu bước ra khỏi bóng tối. Và chúng không đi lang thang vô nghĩa.

Có lẽ\dots\ chúng đang tìm đúng người. Người mà có thể giúp chúng giải thoát được, một lần và mãi mãi.

\begin{center}
  \pov{Hikawa Kuon}
  \uline{Một góc thị trấn Aogami}
\end{center}

Tôi vác hai túi đồ nặng trĩu, khựng bước để hai cô gái đang túm đầu tám chuyện kịp song song phía trước.

Túi trái là thực phẩm cho bữa tối của gia đình tôi: thịt heo thái mỏng để xào gừng, khay cá saba ướp muối, cải xanh, hành lá, đậu hũ trắng mịn và bịch mì udon tươi còn vương phấn bột.

Túi phải thì\dots\ khỏi nói, chứng tích ``đi chợ cùng đám trẻ'': cơm nắm tam giác, khoai tây chiên đủ vị phủ tương, bánh que phủ socola bóng mỡ, vài lon nước ngọt lạnh toát, hộp sữa nhỏ in họa tiết hoạt hình và bịch kẹo dẻo hình ngôi sao Akane-san nằng nặc đòi mua ``để nhìn thôi cũng vui''.

``Này Kuon-kun, tối nay nhà cậu có cá hả?'' Akane-san nghiêng đầu, đuôi mắt nhấp nháy như đèn pin mini.

``Ừ, cá nướng và canh đậu phụ. Đơn giản nhưng đủ no.''

``Yatta\~{}!'' Cô nàng giơ hai tay, váy xoè như pháo hoa mini giữa hẻm vắng.

``Nghe như cậu định cắm trại ở nhà tôi ấy\dots'' tôi thở dài, nụ cười méo xệch.

Yuka-san khẽ liếc vào túi snack, khẽ cười: ``Mua thế này, khéo ăn hết trước bữa tối.''

``Toàn cho mấy vị khách không mời mà đến,'' tôi đáp lướt, ``chứ tôi cần gì.''

Cô bạn tai chó che miệng cười khúc khích; nàng bờm mèo lắc đầu như kẻ từng trải.

Con hẻm về khu dân cư lặng ngắt, thoang thoảng mùi đất ẩm sau cơn mưa xuân buổi sáng, tỏa mùi ngai ngái như lá non vừa bị giày đạt.

Và rồi\dots

Tôi chợt chậm bước.

Ở góc hẻm phía trước, ngay chỗ bức tường loang lổ rêu xanh, có một bóng dáng nhỏ bé.

Một đứa trẻ. Đúng hơn, một bóng người mang dáng dấp đứa trẻ.

Nó đứng quay lưng lại với bọn tôi, thân hình gầy gò, đầu hơi cúi xuống. Một mái tóc được cột thành hai bím tết, che cả đôi mắt, không lay động dù gió nhẹ lùa qua.

``\dots?''

Tôi chưa kịp lên tiếng thì Yuka-san cũng khựng lại.

Akane-san, vẫn vô tư, hỏi: ``Ơ? Có con nít đứng đó kìa?''

Ánh mắt tôi lập tức trượt xuống mặt đất.

Trên nền đất ẩm, ngay trước góc tường, những dấu chân nhỏ xíu in hằn rõ mồn một. Từng dấu sâu hơn dấu chân người lớn. Như thể ai đó vừa đứng đó rất lâu.

Tim tôi đập mạnh hơn một nhịp.

``Akane-san,'' tôi hạ giọng, ``đứng sau tôi.''

Không khí bỗng lạnh đi rất khẽ. Tựa như một cơn gió lạnh thổi qua, đánh tan tiết trời ấm áp mùa xuân, nhường chỗ cho sự rợn người.

Bóng đứa trẻ kia\dots\ không hề quay lại.

Nó đứng bất động, như một hình cắt bằng bóng tối.

Yuka-san nuốt khan. Rồi cô ấy cất tiếng, chậm rãi, trầm hơn bình thường: ``Ở nơi từng là mảnh đất chết\dots\ Có thể còn thứ gì\dots\ chưa được chôn cất đúng nghĩa?''

Câu nói rơi xuống con hẻm như một viên đá nặng.

Cô nàng tóc nâu chớp mắt vài lần. ``Ý bà là\dots\ ma hả?''

Cậu ấy nói rất nhỏ, không còn giọng đùa nữa.

Tôi siết chặt quai túi đồ.

Trong đầu tôi vụt hiện lên những con phố hoang tàn của Aogami. Những linh hồn lảng vảng. Những dấu vết chưa từng được siêu thoát.

``Có khả năng là vậy,'' tôi đáp ngắn gọn.

Ngay lúc đó\dots

Bóng đứa trẻ dần mờ đi.

Không tan biến kiểu khói. Không biến mất ngay lập tức.

Nó như bị ánh sáng xóa dần, từng chút một, cho đến khi góc hẻm chỉ còn lại bức tường ẩm ướt và\dots\ những dấu chân trống rỗng.

Nhưng chừng đó cũng đủ để có thể đưa ra một thông điệp nào đó, rằng\dots

\begin{center}
  ``Bọn tôi vẫn ở đây. Mong chờ được an nghỉ đúng nghĩa.''
\end{center}

Akane-san nắm chặt tay áo tôi, giọng run run: ``Kuon-kun\dots\ mình vừa thấy cái gì vậy\dots?''

Yuka-san thở dài, ánh mắt trầm xuống. ``Quả thật\dots\ chúng bắt đầu xuất hiện ở khắp mọi nơi rồi. Không chỉ ở trong trường hay ở cánh rừng, mà ở những nơi bình thường này\dots''

Tôi gật đầu. ``Chúng không bị trói buộc bởi một chỗ nữa. Chúng bắt đầu xuất hiện ở nơi từng có ký ức, từng có mất mát.''

Một khoảng im lặng kéo dài. Tiếng xe ngoài đường lớn vang vọng xa xa, như thuộc về một thế giới khác.

Cuối cùng Akane-san lên tiếng, giọng yếu ớt: ``Vậy\dots\ chúng ta vẫn ăn cá nướng tối nay chứ?''

Câu hỏi ngây ngô làm tôi bật cười khẽ, dù trong lòng chẳng nhẹ đi bao nhiêu. ``Ăn chứ, cá thì vẫn phải ăn.'' Tôi xoa đầu cô ấy.

Yuka-san mỉm cười rất nhẹ, nhưng ánh mắt vẫn còn lo lắng.

Tâm trạng của chúng tôi lúc này không còn là sự vô tư của một buổi đi chợ. Đó là cảm giác thế giới đang dần để lộ những vết nứt, và những gì bị chôn giấu đang bắt đầu bước ra ánh sáng.

\ct

``\textbf{Trời ơi dễ thương quá!}'' tiếng reo của Kana-san như tia lửa đốt cháy không khí. Cô lao thẳng vào kệ thú bông, váy xoè như cánh hoa rụng, đến mức rơi cả mũ lưỡi trai hình quỷ trên đầu.

Ngay sau đó, Kanade-san giơ cao chùm móc khóa lấp lánh, mắt cười híp lại: ``Cái này treo cặp sách là hợp cả!''

Cửa hàng lưu niệm nép mình bên hông trung tâm thương mại, nhỏ xíu nhưng rực rỡ. Kệ gỗ sáng bóng chất đầy móc khóa hình thú, thú nhồi bông, tượng chibi, bưu thiếp phong cảnh; chuông gió lúc lắc reo leng keng mỗi khi cửa xoay mở. Vừa bước qua ngưỡng, cả nhóm tản ra như hạt cát dưới nắng. ``Cặp thiên sứ - ác ma'' (như Onii-sama vẫn thường gọi) chạm mặt đã phấn khích ôm trọn giá móc khóa.

Yuki-san xoay những cầu tuyết tí hon, tuyết giả bay lả tả trong ánh đèn. Còn Mitsuki-san đứng cạnh tôi, khoanh tay, mắt lặng lẽ săm soi mọi thứ, khóe môi khe khẽ nhếch.

``Trông bọn họ không khác gì con nít nhỉ,' tôi bật cười nhẹ bình luận.

``Mấy cậu đó cứ nhìn thấy gì đó dễ thương là lại réo lên thôi na\~{}'' cô bạn tóc hồng ngân nga đáp lại, hướng mặt về phía Hotaru-san đang nhìn ngắm những con thú bông với đôi mắt sáng rỡ. ``Đặc biệt là cái tên Hotaru kia kìa.''

Tôi không đi đâu xa cả. Chỉ dựa vào kệ, hít hà mùi gỗ mới và vải bông thơm tho, để tiếng cười nói hoà nhạc nền du dương quện vào nhau thành một tấm mền ấm phủ xuống vai.

Cô nàng da ngăm mà chúng tôi nhắc tới ghé góc tủ, nhấc lên bức tượng chibi áo choàng cũ kỹ, mắt to tròn long lanh. Cô ngắm nghía một lúc, bật cười khẽ: ``Nhân vật này\dots\ làm tôi nhớ hồi bé ghê.''

Kana-san quay phắt lại: ``Ngày xưa gì cơ?''

Hotaru-san nghiêng đầu, giọng mơ hồ như kể một giấc mộng: ``Hồi thích mấy anh hùng cô đơn, lúc nào cũng tưởng mình đặc biệt.''

Kanade vỗ đùng: ``Tuổi thơ của cậu dữ dội lắm nhỉ?'' Cả bọn cười rộ, tiếng vang lắc lư trên xà nhựa.

Tôi liếc Mitsuki-san. Cô ấy mỉm cười kín đáo của một người đứng ngoài nhưng vẫn cảm được niềm vui.

``Nhóm này ồn ghê,'' tôi thì thầm.

``Nhưng vui,'' cô đáp, mắt sáng lên. ``Lâu lắm mới thấy thế.''

Tôi gật nhẹ. Nếu không có những điều kỳ lạ từ khi đặt chân tới thị trấn này, buổi chiều hôm nay đã có thể tròn trịa như một hạt ngọc trai, bình thường, đời thường, và hoàn hảo.

\dots Chỉ tiếc rằng, mọi thứ lại không diễn ra suôn sẻ như vậy.

*CỘP*

Một âm thanh khô khốc vang lên, không lớn nhưng đủ rõ để cắt ngang tiếng cười nói.

Một con búp bê vải mặc váy xanh nhạt, tóc nâu thắt hai bên bất ngờ nằm lăn trên sàn gỗ, rơi từ kệ cao ngang ngực người lớn.

Hotaru-san đang đứng cách cái kệ đó gần hai mét, Yuki-san thì ở phía bên kia cửa hàng, còn tôi và Mitsuki-san chưa hề bước tới gần. Tôi nhìn lên kệ: những con búp bê còn lại vẫn xếp ngay ngắn, không dấu hiệu bị xô lệch, không ai vừa chạm vào. Một cơn lạnh mỏng manh chạy dọc sống lưng tôi.

Nàng tóc tím giật mình: ``Ơ? Ai làm rơi vậy?''

Cô bạn thân của cô ấy quay phắt lại: ``Có ai đụng đâu?''

Cô bạn tóc đen để tay lên cằm, suy đoán: ``Chắc là có ai đó xếp không ngay ngắn thôi\~{}''

Cô nàng tóc hồng khẽ thì thầm: ``\dots Tự rơi xuống ư?''

Tôi không trả lời, chỉ nhìn chằm chằm vào con búp bê nằm trên sàn. Nó nghiêng đầu sang một bên, đôi mắt khâu chỉ đen cong cong trông như đang cười. Tôi linh cảm rằng nó không rơi do sơ suất, mà bị một thế lực nào đó thả xuống.

Không khí trong cửa hàng vẫn râm ran, nhưng với tôi, mọi âm thanh bỗng dưng đổi khác một cách chóng mặt. Một cái gì đó lạnh mỏng, nhẹ, như sương sớm tràn qua mắt cá chân.

Tôi ngẩng đầu lên.

Ở khoảng trống giữa hai dãy kệ gỗ, nơi ánh đèn vàng chiếu không tới hết, một bóng mờ như sương đang đứng đó. Nó không rõ hình khối hoàn toàn, cũng không đặc như cái bóng ở sân trường, lại càng không rõ ràng như đứa trẻ ở góc ngõ mà chúng tôi từng chạm trán. Nó giống như một lớp khói mỏng tụ lại thành hình người. Cái bóng thấp ngang eo người lớn, dáng nhỏ, khoác thứ gì đó trông như áo vải rộng thời rất xưa. Phần tà mờ loang vào không khí, như bị gió vô hình xé nhẹ. Không có tiếng thở, không chuyển động, chỉ đứng.

Thứ khiến tim tôi thắt lại chính là khuôn mặt. Đôi mắt của em ấy bị che mất bởi làn tóc mái dài. Kỳ lạ thay, tôi vẫn biết nó đang quan sát, không hề mang ác ý, không hăm dọa, giống như một linh hồn đứng lặng nhìn người sống trôi qua.

Tôi nghe rõ tim mình hụt một nhịp. Bàn tay cầm con búp bê run lên. Hotaru-san và mấy người kia vẫn đang cười nói phía sau, họ chưa nhìn thấy gì cả. Một phần trong tôi gào lên rằng phải chạy, nhưng chân tôi không nhúc nhích.

Thay vào đó, tôi hít thật khẽ, ép môi cong lên thành một nụ cười run rẩy, rồi thì thầm rất nhỏ, như sợ làm tan lớp sương ấy: ``\hl{Xin lỗi\dots{} vì đã làm phiền.}''

Bóng mờ không phản ứng. Tôi nuốt khan: ``\hl{Nếu em muốn nói gì\dots{} chị sẽ lắng nghe.}''

Không khí lặng như tờ. Tôi tiếp tục, giọng khẽ đến mức chính mình cũng suýt không nghe rõ: ``\hl{Không biết em có\dots{} giống những đứa trẻ mà Sakura-san từng nhắc đến không?}''

Trong khoảnh khắc đó, lớp sương quanh hình dáng ấy khẽ rung lên, giống như mặt nước bị một giọt mưa rơi xuống. Rồi, cái bóng mờ dần, tan vào ánh đèn trong cửa hàng, như chưa từng tồn tại.

Tôi đứng chết lặng vài giây, tim vẫn đập loạn xạ. Trong lòng không có cảm giác nguy hiểm, chỉ là một nỗi buồn sâu thẳm không thành lời.

``''\textbf{Suzuki-san! Lại đây xem này!}''

Tiếng Yuki-san kéo tôi về thực tại. Tôi hít sâu, đặt con búp bê lên kệ, lần này cẩn thận hơn, rồi bước về phía mọi người.

Cuộc ``phiêu lưu'' lại tiếp tục. Cả nhóm kéo sang cửa hàng phụ kiện kế bên, rồi tiệm kem, rồi loanh quanh khu gắp thú. Kana-san la hét mỗi lần gắp hụt, Kanade-san cổ vũ ầm ĩ như đang xem chung kết, Hotaru-san cười đến chảy nước mắt mỗi lần ai đó nổi cáu.

Bề ngoài mọi thứ vẫn bình thường, nhưng trong đầu tôi hình ảnh bóng mờ cứ lặp đi lặp lại. Những linh hồn không tấn công, không đe dọa, chỉ lặng lẽ quan sát như những ký ức chưa được yên nghỉ.

Tôi nhớ những câu chuyện cũ về hủ tục, về đứa trẻ bị chôn sống trong nghi lễ cổ xưa, về những linh hồn không ai gọi tên. Phải chăng những gì bọn tôi gặp chính là họ, những nạn nhân bị bỏ lại giữa lịch sử và lãng quên?

Tôi bước theo bạn bè, cười nói như không có gì, nhưng trong lòng câu hỏi vẫn treo lơ lửng: nếu họ vẫn quanh quẩn, điều gì mới có thể giúp họ thật sự yên nghỉ? Và liệu bọn tôi có đang vô tình bước sâu hơn vào thế giới của những linh hồn ấy?

\begin{center}
  \pov{-}
  \uline{Năm giờ chiều.}
\end{center}

Ở nơi tập kết ban sáng, nhóm của Nanase đứng thành một vòng tròn nhỏ. Không khí vẫn vui vẻ, nhưng đã lắng hơn lúc ban trưa.

Honoka vươn vai thật mạnh. ``Trời ơi, xem phim xong rồi đi chơi nguyên buổi chiều luôn. Đúng kiểu ngày nghỉ mà tớ muốn đó!''

Shiina khoanh tay, liếc về phía rạp. ``\dots Ừ. Nếu bỏ qua cái cảnh kỳ quặc kia.''

Nanase gật nhẹ, nét mặt nghiêm hơn thường lệ. ``Cảm giác\dots\ như có thứ gì đó thật sự ở đó. Một cái gì đó làm tôi rợn tóc gáy ấy.''

Shino đứng cạnh Kaigou, hai tay đan trước ngực. Nụ cười của cô nàngg vẫn dịu dàng, nhưng ánh mắt thoáng trầm xuống. ``Mình nghĩ đó cũng là một trải nghiệm đáng nhớ đấy chứ. Dù gì đây cũng là lần đầu ta đi chơi với nhau mà.''

Kaigou hừ nhẹ, tay đút túi áo khoác. ``Dù sao thì ta vẫn chưa hẳn quen nhau mà.
Nhưng mà nhá\dots''

Cô dừng một nhịp. ``\dots chuyến đi hôm nay vui thật.''

Nàng Cáo sa mạc cười toe: ``Thấy chưa! Vừa vui vừa rợn mới đáng nhớ chứ!''

Trong khi đó, nàng Sư tử liếc cô ấy. ``Bớt hào hứng với mấy thứ siêu nhiên giùm cái.''

Mọi người bật cười. Rồi, đến lúc chia tay. Mỗi người bắt đầu về một hướng.

Kaigou quay sang cô nàng tóc nâu.
``Tôi sẽ về hướng này.''

Shino gật đầu. ``Mình cũng vậy. Hôm nay\dots\ cảm ơn cậu.''

Cô nàng tóc đuôi ngựa xoa nhẹ đầu cô rất nhanh, như sợ bị phát hiện là dịu dàng quá mức. ``Đi cẩn thận.''

``Ừm. Hẹn gặp cậu vào thứ hai.''

Hai đứa rẽ về hai con đường khác nhau.

Nhưng dù bước xa dần, cả hai vẫn không thể hoàn toàn gạt khỏi đầu hình ảnh bóng trẻ nhỏ co gối trong rạp chiếu phim.

Một thứ gì đó chưa chịu rời đi.

\ct

Ở phía bên kia khu trung tâm thương mại, nhóm Suzuki cũng chuẩn bị ra về.

Hotaru khoác tay cô thật tự nhiên: ``Thôi, Onee-sama giao nhiệm vụ rồi nha. Hôm nay tôi hộ tống cậu về tận nơi.''

Suzuki khẽ cười: ``Làm như mình là trẻ con không bằng\dots''

``Không phải,'' cô nàng da ngăm cười híp mắt, ``bảo vệ báu vật của lớp đấy. Trông cậu dễ thương thế này cơ mà.''

Kana huýt sáo trêu: ``Ghê chưa, được ưu ái dữ ha!''

Kanade gật gù phụ họa: ``Hotaru lúc nào cũng ngầu theo kiểu kỳ lạ như thế.''

Mitsuki vẫy tay dịu dàng: ``Về cẩn thận nhé, Suzuki.''

Yuki đã bắt đầu bàn kế hoạch karaoke lần sau, nên chỉ vẫy tay chào tạm biệt.

Trên đường về, Hotaru nói không ngừng nghỉ: mấy món ăn ngon, cú gắp thú hụt thảm họa của Kana, kế hoạch lần sau nhất định phải kéo cả lớp đi chơi tiếp. Trái ngược với đó, Suzuki chỉ ``ừm'', ``vâng'', ``thật à'' cho có lệ, tâm trí vẫn kẹt lại trong cửa hàng lưu niệm, nơi bóng mờ như sương đứng lặng, không mắt, không chuyển động, chỉ lặng lẽ quan sát: một linh hồn không làm hại ai nhưng mang theo nỗi buồn cổ xưa.

Hotaru đang thao thao bỗng khựng lại: ``Ê, Suzuki.''

``Hả?''

``Cậu ổn không? Nãy giờ thấy cậu cứ lơ đãng sao ấy.''

Cô nàng mắt hai màu giật mình, mỉm cười nhẹ: ``Ừ\dots\ chỉ là đang nghĩ mấy chuyện kỳ lạ thôi.''

Hotaru nghiêng đầu, không hiểu tôi nói gì. Suzuki không giải thích thêm, chỉ cười mỏng. Nhưng sâu trong thâm tâm, cô biết rất rõ, đúng là vậy.

\ct

Ở nhà Hikawa, không khí đã khác hẳn.

Kuon đứng khoanh tay trước cửa phòng khách, nói: ``Rồi, vui chơi đủ rồi. Về nhà đi kẻo bố mẹ lo.''

Akane ngồi bệt xuống sàn như dính keo. ``Không về đâu\~{} Bố mẹ tôi hôm nay không về\~{}''

Yuka khoanh tay gật gù: ``Akane không về thì tôi cũng không về đâu. Tôi cũng muốn ăn thử món cá mà Kuon-kun làm nữa\~{}''

Miku đang dọn đống gói bánh trên bàn, đáp lại với giọng ấm: ``Ở thêm chút cũng chẳng sao đâu. Tôi sống một mình mà.''

Inari đứng sau, nhỏ giọng: ``Nếu\dots\ mọi người ở lại thì mình cũng vậy\dots''

Cậu Tổng nhướng mày, nhìn cô nàng bờm Cáo trắng với sự ngạc nhiên: ``Tưởng bố mẹ cậu nghiêm khắc lắm chứ? Tôi nghe cậu đã từng nói vậy mà?''

``Ừ thì\dots'' Inari đáp lại, ``nhưng mà\dots\ mình đã xin phép cậu sang nhà để học nhóm rồi. Bố mẹ mình cũng không phản đối gì.''

``Nói dối tốt thật đấy,'' cậu Tổng đáp lại, thở dài. ``Vậy mấy cô nưong còn lại thì sao đây?''

Trái ngược lại với bốn người bờm thú, Hanabi đã xỏ giày từ sớm: ``Thôi, tôi về vậy. Không muốn chết đói nếu Kuon-kun nổi cáu.''

Kaoru theo sau: ``Gặp lại tuần sau nhé.''

Cậu Tổng không đáp gì, chỉ ra hiệu kiểu như ``lượn đi cho nước trong.'' Cửa đóng sập, cái *CẠCH* khô khốc.

Trong phòng chỉ còn bốn người mà Kuon thầm gọi là ``GAMERS ERROR''. Kuon chống hông, cười méo xệch: ``Các cậu định cắm trại ở nhà tôi thật luôn à?''

Nàng bờm Chó lăn qua lăn lại trên sàn: ``Ở tới tối đi mà.''

Nàng bờm Mèo hùa theo, lăn qua lăn lại với cô bạn: ``Ăn tối xong sẽ về. Hứa.''

Nàng bờm Sói bình thản: ``Chúng tôi sẽ không gây rối đâu.''

Cuối cùng, nàng bờm Cáo trắng run run giơ tay: ``Để\dots\ để mình rửa bát cho.''

Kuon im lặng vài giây, nhìn bốn khuôn mặt bám dai như đỉa, rồi xoa trán chịu thua: ``Ăn tối xong phải về. Không thương lượng.''

Akane bật dậy giơ tay: ``\textbf{Hoan hô!}'' Yuka đập tay cô cái bốp. Miku thở phào. Inari cười rụt rè như mèo lọt cửa.

Kuon lẩm bẩm: ``\textit{Tự dưng ngày nghỉ mà phải phục vụ bốn con giời làm gì không biết\dots}'' Nhưng, khóe môi cậu khẽ cong, khe khẽ, ``\textit{làm mình liên tưởng tới hồi nhóm họ tới nhà ngủ hồi dịp Tết năm ngoái.}''

Buổi chiều dần trôi sang tối.

Ánh nắng cuối ngày đổ dài qua khung cửa sổ nhà Hikawa, nhuộm căn phòng khách thành một màu cam nhạt yên bình, trái ngược hoàn toàn với những gì đang âm thầm tích tụ trong tâm trí ba anh em.

\ct

Khi kim đồng hồ gần chạm sáu, Kaigou và Suzuki lần lượt bước vào nhà, không réo rắt trêu chọc như mọi khi, chỉ gật đầu nhẹ như ra ám hiệu, rồi lặng lẽ đi thẳng vào phòng Kuon, nơi vẫn được dùng để bàn những chuyện không dành cho người ngoài nghe. Cánh cửa khép hờ, không khí nặng hẳn, nghiêm trang hơn ban ngày rất nhiều.

Kuon dựa lưng vào thành giường, Kaigou khoanh tay dựa tường, Suzuki ngồi xếp bằng trên thảm, hai tay đặt lên đầu gối. Im lặng kéo dài một nhịp thở, rồi cậu con trai cất tiếng trước: ``\vie{Hai đứa hôm nay thế nào?}''

Rõ ràng, điều mà cậu nhắc tới không hoàn toàn là buổi đi chơi, mà là những \emph{hiện tượng kỳ quặc} mà cả ba người họ bắt gặp được. Cậu dám chắc Kaigou và Suzuki cũng muốn biết.

Suzuki kể trước, giọng nhỏ như gió lọt khe cửa: về trung tâm thương mại ồn ào tiếng cười, về con búp bê bất ngờ rơi khỏi kệ chẳng ai động vào, và về bóng mờ như sương không mắt nhưng vẫn khiến người ta thấy bị nhìn chằm chằm. ``\vie{Hồn ma mà em gặp đó\dots\ không có mắt,}'' cô thì thầm, ``\vie{nhưng em có cảm giác nó đang nhìn mình.}'' Cô mô tả thêm: bóng hình thấp thoáng như khói, mờ đến mức nếu không để ý sẽ tưởng chỉ là bóng đèn rung nhẹ, nhưng khi nhìn kỹ lại thấy nó có dáng một đứa trẻ, tay chân mảnh dẻ, khoác áo choàng rộng thùng thình, tóc xõa dài che khuất nửa khuôn mặt — hay đúng hơn, nơi lẽ ra phải có mặt ấy chỉ là một mảng sương trắng mờ, gương mặt trống trơn, chỉ là một khoảng trống nhưng vẫn khiến người ta biết rằng nó đang nhìn. Nhưng không phải nhìn bằng ánh mắt, mà như thể nó đang nhìn xuyên vào ký ức, vào nỗi buồn mà người sống cố giấu.

Chẳng cần cô nói thẳng, cả hai anh chị lớn đều hiểu ám chỉ cô em út nói tới là gì.

Kaigou khẽ tặc lưỡi, đổi lượt: ``\vie{Tôi thì đối mặt nó ở rạp chiếu phim.}'' Cô kể hàng ghế cuối trống không, bóng tối đậm dần khi màn hình đổi cảnh, và hình ảnh đứa trẻ co gối hiện ra thình lình giữa bộ phim tình cảm học đường chẳng liên quan gì đến nỗi buồn ấy. ``\vie{Shino-san lúc đó không hoảng,}'' cô nói khẽ, ``\vie{như thể đã đoán định được rồi.}'' Cô nhớ lại ánh mắt của Shino: không sợ hãi, không hoảng loạn, chỉ là một cái nhìn sâu thẳm như hồ nước tối, như thể cô ấy đã từng gặp hồn ma này, hay ít nhất là đã từng nghe tiếng nó kêu gọi trong giấc mơ. Shino không nhìn thẳng vào bóng ma, nhưng vai cô rung lên một cái thật nhẹ, như gió lạnh lướt qua, và đôi mắt cô ánh lên một nỗi đau quen thuộc của cả một thời đại đã chôn vùi.

Kuon siết nhẹ tay, đến phiên cậu: ``\vie{Tôi với Yuka-san và Akane-san thấy một đứa trẻ đứng quay lưng ở góc ngõ,}'' cậu kể tiếp, giọng trầm xuống khi nhắc đến những dấu chân nhỏ in hằn trên đất ẩm. Cậu mô tả: bóng hình đứng im như tượng, không quay lại, không cử động, chỉ lặng lẽ hiện diện như một vết thương cũ của con hẻm. Yuka là người đầu tiên nhận ra điều không ổn. Cô không la hét, chỉ bước chậm lại, mắt dán vào bóng trẻ, môi mấp máy như thể đang đọc một cái tên không còn ai nhớ. Rồi cô nói nhỏ, gần như tự lẩm bẩm: ``\vie{Nó không có mặt\dots\ nhưng nó vẫn đang khóc.}'' Kuon không nghe thấy tiếng khóc, nhưng cậu cảm được một làn gió lạnh thổi qua từ những nấm mồ không bia đá.

Căn phòng chìm trong tĩnh lặng. Ba câu chuyện khác biệt, nhưng điểm chung hiển hiện: đều là trẻ nhỏ, đều không tấn công, đều xuất hiện gần những nơi từng mắc kẹt trong vòng lặp. Suzuki lại thì thầm: ``\vie{Như thể\dots\ họ đang tìm ai đó.}'' Kaigou gật chậm: ``\vie{Hoặc muốn nói điều gì đó.}''

Kuon nhìn xuống sàn: ``\vie{Và tất cả đều gắn với những nơi từng xảy ra chuyện.}''

Không ai thốt ra, nhưng cả ba cùng nghĩ tới những hủ tục cũ, những đứa trẻ bị chôn sống, những mảnh đất chết chưa từng được yên nghỉ. Chúng dường như cố gắng muốn thổ lộ một điều gì đó, như đang ra dấu hiệu cầu cứu.

Không khí đặc quánh lại; những bí ẩn hôm nay không rời rạc nữa\dots

Chúng đang kết nối với nhau, như một sợi dây liên kết kỳ lạ.

*BZZT*

Tiếng rung vang lên từ cổ tay Kuon. Là chiếc đồng hồ thông minh cậu vẫn đeo, nhưng lần này, màn hình không hiện giờ, mà giờ cũng cất giọng: ``\eng{Có vẻ như những linh hồn bắt đầu trà trộn vào chúng ta rồi.}''

Ba người giật mình. Kuon giơ tay lên. Giọng nói trầm đặc pha chút hơi Pháp tiếp tục: ``\eng{Dường như chúng đang tìm cách nói với chúng ta điều gì đó. Ít nhất, theo hồn ma mà Kuon-user chạm trán.''}

Suzuki nuốt khan, giọng run run: ``\eng{Ý\dots\ Ý ông là sao?}''

Kuon chau mày, mắt dán vào mặt đồng hồ. ``\eng{Game-san\dots\ ông hiểu được bọn chúng nói gì à?}''

Màn hình loé sáng, một dòng chữ nhỏ xíu hiện rồi mờ dần: đi vào rừng. Sau đó, giọng nói vang lên không trầm hơn, như thể đọc lại một bản tin cũ: ``\eng{Họ muốn các cậu đi vào trong rừng. Nhưng chúng cũng không nói rõ là định làm gì.}'' Ông tiếp tục: ``\eng{Các cậu đang ở giữa một mạng lưới những linh hồn bị bỏ rơi. Và họ đã bắt đầu biết cách tìm đường về phía các cậu.}''

Im lặng.

Rồi màn hình tắt ngóm, như vừa rút lui vào bóng tối.

Kaigou thì thầm, giọng khàn: ``\vie{Nếu như những gì ông ta nói là sự thật\dots\ Vậy thì\dots}''

Kuon không đáp. Cậu chỉ nhìn vào màn hình đã tối, như thể trong đó vẫn còn một đôi mắt đang nhìn lại.

Và ngoài cửa sổ, gió bỗng thổi mạnh hơn, đập vào khung kính như tiếng gõ cửa.

Họ đang đến gần hơn tới một sự thật nào đó vẫn còn đang che giấu, một điều gì đó vẫn còn uẩn khúc đằng sau những lời kêu cứu của những đứa trẻ hồn ma đó.

\dots

\dots

Tiếng Akane réo vang từ tầng dưới, xé toang bầu không khí nặng trịch: ``\textbf{Kuoooooon-kuuuun!! Bao giờ ăn tối vậy!}''

Giọng Yuka tiếp ngay sau, nhão nhoẹt: ``\textbf{Tôi đói rồi nhaaa\~{}}''

Miku gõ cửa cộc cộc, nhịp đều đặn: ``Đồ ăn nguội sẽ không tốt cho sức khỏe đâu đó.''

Inari thì thào theo sau, rụt rè: ``M-Mình phụ dọn bàn được không\dots?''

Không khí vừa căng như dây đàn lập tức xì hơi, trở lại nhịp sống thường nhật.

Suzuki bật cười khẽ, vai rung rung.

Kaigou thở dài, cau mày: ``Đúng là không cho người ta nghiêm túc lâu được.''

Kuon đứng dậy, vỗ nhẹ tay vào đầu gối: ``Thôi. Ra ăn trước đã.'' Cậu khựng lại, giọng trầm hơn: ``Chuyện kia để sau.''

Cánh cửa phòng mở toang, ánh đèn vàng ấm tràn vào, kéo theo mùi cơm nóng và tiếng ồn ào quen thuộc.

Nhưng dù bữa tối đã bày ra đầy đủ, ba người kia vẫn biết rõ: những linh hồn kia chưa hề biến mất; chúng chỉ đang chờ.

Đêm buông xuống thị trấn Aogami, yên tĩnh đến bất thường. Ở đâu đó giữa những con hẻm vắng, dưới ánh màn ảnh đã tắt, hay trong góc kệ cửa hàng cũ kỹ, những cái bóng nhỏ bé vẫn lặng lẽ quan sát. Không hận thù, không tấn công, cũng không có thù địch gì, chỉ mang theo những câu chuyện chưa ai kể và những cái chết chưa từng được lắng nghe.

Ngày mai, bí ẩn sẽ không còn đứng yên nữa.
\end{document}